\chapter{Antilinear Representations: The Case of $\bC / \bR$}\label{chapter:tenfoldway}

In this chapter we continue with the setup of Chapter \ref{chapter:twelvefoldway}, and specialise to the case when $\bL / \bK$ is the extension $\bC / \bR$. In this case, we refer to semilinear representations as \emph{antilinear} representations.

\begin{itemize}
    \item 
In Section \ref{sect:antilinearoverview} we recall the key properties that antilinear representations satisfy as opposed to general semilinear representations.

\item In Section \ref{sect:recoverDyson10} we compare Dyson's Tenfold Way and our classification into types.

\item In Section \ref{sect:antilinearvsrealreps} we consider results for $\bR$-linear representations and use Dyson's Tenfold Way to understand how these extend to results for general antilinear representations.

\item In Section \ref{sect:realisablecyclo}, we prove an induction theorem for antilinear representations and use this to show that antilinear representations of finite groups admit models over cyclotomic fields.

\item In Section \ref{sect:unitaryreps} we extend the theory of the Frobenius-Schur indicator to unitary representations of certain infinite groups, including compact groups and discrete amenable groups.

\item In Section \ref{sect:deponrealstr} we study for a fixed group $H$ how the D-Type of a representation of $H$ varies as we alter the Real structure on $H$.
\end{itemize}

\section{Overview of Antilinear Representations}\label{sect:antilinearoverview}

The two main facts which restrict the structure of antilinear representations as compared to general semilinear representations for degree two extensions are that

\begin{itemize}
    \item $\bC$ is algebraically closed, so any $V \in \Irr_{\bC}(H)$ is absolutely irreducible,
    \item Any finite-dimensional division algebra over $\bR$ is either $\bR$, $\bC$ or $\bH$ (Frobenius' Theorem).
\end{itemize}

From the absolute irreducibility, the possible types that can occur are a subset of those twelve types of Definition \ref{def:twelvetypes}, and so for antilinear representations we use the D-type notation instead of the type. It actually happens that D-types II.V and IX.V cannot occur (and the remaining ten D-types do), but this fact, known as Dyson's Tenfold Way, is not immediate, and relies of the classification of $C_2$-graded division algebras over $\bR$ (see Section \ref{sect:recoverDyson10} below).

From the classification of finite-dimensional division algebras over $\bR$, the three possible S-types for an antilinear representation $W$ are characterised by the endomorphism ring of $W$ being $\bR$, $\bC$ or $\bH$, and in particular the division algebras appearing in Table \ref{tab:24types6} are either $\bR$, $\bC$ or $\bH$.

The other main feature of the case when $\bL / \bK$ is $\bC / \bR$ is the existence of the Frobenius-Schur indicator. If $H$ is finite, and $V \in \Irr_{\bC}(H)$ has character $\chi$, we define
\[
\sF^{\sharp}(V) \coloneqq \sF^{\sharp}(\chi) \coloneqq \sum_{g \in G \setminus H} \chi(g^2).
\]

This invariant allows one to compute the S-type (cf.\ Theorem \ref{thm:threeoptions}) of any $V \in \Irr_{\bC}(H)$.

\begin{prop}{\cite[Thm.\ 4.2]{RUMTAY}}\label{prop:FSfinitegroups}
Suppose that $H$ is finite and $V \in \Irr_{\bC}(H)$. Then:
\begin{itemize}
    \item $\fS(V) = \mathbf{A}$ if and only if $\sF^{\sharp}(V) = 1$,
    \item $\fS(V) = \mathbf{B}$ if and only if $\sF^{\sharp}(V) = 0$,
    \item $\fS(V) = \mathbf{C}$ if and only if $\sF^{\sharp}(V) = -1$.
\end{itemize}
\end{prop}

As a consequence of the induction theorem, we prove in Section \ref{sect:realisablecyclo}, we obtain that Proposition \ref{prop:FSfinitegroups} also applies when $H$ is finite to certain cyclotomic extensions (Theorem \ref{thm:15waysforcyclotomic}, Remark \ref{rem:FSforcyclo}).

We also extend Proposition \ref{prop:FSfinitegroups} to certain infinite groups in Section \ref{sect:unitaryreps}.

\section{The Tenfold Way for $\bC / \bR$}\label{sect:recoverDyson10}

In this section we give a simple proof of the tenfold way for $\bC / \bR$ by showing that for $\bC / \bR$ cases II.V and IX.V of Definition \ref{def:twelvetypes} cannot occur. The key input is the classification of $C_2$-graded division algebras over $\bR$, which we now recall.

By Lemma \ref{lem:centreofC2graddiv}, the graded centre of any $C_2$-graded division algebra over $\bR$ is either $\bR$ or $\bC$, and so any $C_2$-graded division algebra can be viewed as an element of $\BrW(\bR)$ or $\BrW(\bC)$ (see Definition \ref{def:BW}). As with the ordinary Brauer group, any class of $\BrW(\bF)$ has a unique representative by a central graded division algebra over $\bF$ by Corollary \ref{cor:divalgrep}, and therefore the set of $C_2$-graded division algebras over $\bR$ is canonically identified with the disjoint union
\[
\BrW(\bR) \sqcup \BrW(\bC).
\]

From the exact sequence \eqref{eqn:exactseqbrauer} of Appendix \ref{appendix:gradedalgebras}, we have that $\BrW(\bC) \cong \bZ / 2 \bZ$, and that $\BrW(\bR)$ fits into an exact sequence
\begin{equation}\label{eqn:exactBrWR}
    0 \rightarrow \bZ / 2 \bZ \rightarrow \BrW(\bR) \rightarrow \bZ / 4 \bZ \rightarrow 0,
\end{equation}
as $\Br(\bR) \cong \bZ / 2 \bZ$ and $Q(\bR)$ is cyclic with generator $(1,1)$.

To produce non-trivially graded elements of $\BrW(\bR)$ and $\BrW(\bC)$, we consider the following.

\begin{defn}
For an finite-dimensional division algebra $A$ over $\bR$, we set
\[
A[\epsilon] \coloneqq A[\epsilon] / (\epsilon^2 - 1), \qquad A[\delta] \coloneqq A[\delta] / (\delta^2 + 1),
\]
where $\epsilon$ and $\delta$ are variables which commute with all elements of $A$. We view these as a $C_2$-graded $\bR$-algebras where $A$ has degree $0$ and $\epsilon$ and $\delta$ have degree $1$.

When $A = \bC$, we also write 
\[
\bC_{\hspace{-0.1em}\rtimes \!}[\epsilon], \qquad \bC_{\hspace{-0.1em}\rtimes \!}[\delta] 
\]
for the same abelian groups, but now viewed as $\bR$-algebras via the skew multiplication:
\[
\epsilon z = \overline{z} \epsilon, \: \epsilon^2 = 1, \qquad \qquad \delta z = \overline{z} \delta, \: \delta^2 = -1.
\]
We view $\bC_{\hspace{-0.1em}\rtimes \!}[\epsilon]$ and $\bC_{\hspace{-0.1em}\rtimes \!}[\delta]$ as $C_2$-graded $\bR$-algebras where $\bC$ has degree $0$ and $\epsilon$ and $\delta$ have degree $1$.
\end{defn}

We note that $\widehat{Z}(A[\epsilon]) = Z(A) = \widehat{Z}(A[\delta])$, and $\widehat{Z}(\bC_{\hspace{-0.1em}\rtimes \!}[\epsilon]) = \bR = \widehat{Z}(\bC_{\hspace{-0.1em}\rtimes \!}[\delta])$.

We can use this to state the classification of $C_2$-graded division algebras over $\bR$. Motivated by Corollary \ref{cor:Dalphaisom}, for such a $D$ we also compute the algebra $D_i$, which we use in Theorem \ref{thm:dyson10} below.

\begin{prop}\label{prop:gradeddivalgoverR}
The $C_2$-graded division algebras over $\bR$ up to isomorphism are listed in the first column of Table \ref{tab:graddivoverR}. The group $\BrW(\bR)$ is generated by $\bR[\epsilon]$, and for $1 \leq k \leq 8$, the $C_2$-graded division algebra over $\bR$ in the $k$th row is equivalent in $\BrW(\bR)$ to the $k$th power of $\bR[\epsilon]$.

For $D$ such an algebra, the graded centre $\widehat{Z}(D)$, $j(D) \in Q(\widehat{Z}(D))$, ungraded part $D_0$, isomorphism class as semisimple $\bR$-algebra $D_\text{ss}$, and associated $\bR$-algebra $D_i$, are as in Table \ref{tab:graddivoverR}. 
\end{prop}

\begin{table}
\caption{D-Types and $C_2$-graded division algebras over $\bR$}
\label{tab:graddivoverR}
\begin{center}
\def\arraystretch{1.1}
\begin{tabular}{ c | c c c c c | c} 
    \toprule
$D$ & $\widehat{Z}(D)$ & $j(D)$ & $D_0$ & $D_{\text{ss}}$ & $D_{i}$ & $\fD(W)$ \\
\midrule
 $(\bR)$ & $\bR$ & $(0,1)$ & $\bR$ & $\bR$ & $\bC$ & V \\ 
 $\bR[\epsilon]$ & $\bR$ & $(1,1)$ & $\bR$ & $\bR \times \bR$ & $M_2(\bR)$ & I \\ 
 $\bC_{\hspace{-0.1em}\rtimes \!}[\epsilon]$ & $\bR$ & $(0,-1)$ & $\bC$  & $M_2(\bR)$ & $M_2(\bR) \times M_2(\bR)$ & III \\ 
 $\bH[\delta]$ & $\bR$ & $(1,-1)$ & $\bH$  & $M_2(\bC)$ & $M_4(\bR)$ & II \\ 
 $(\bH)$ & $\bR$ & $(0,1)$ & $\bH$  & $\bH$ & $M_2(\bC)$ & VI \\ 
 $\bH[\epsilon]$ & $\bR$ & $(1,1)$ & $\bH$  & $\bH \times \bH$ & $M_2(\bH)$ & VIII \\ 
 $\bC_{\hspace{-0.1em}\rtimes \!}[\delta]$ & $\bR$ & $(0,-1)$ & $\bC$  & $\bH$ & $\bH \times \bH$ & X \\ 
 $\bR[\delta]$ & $\bR$ & $(1,-1)$ & $\bR$  & $\bC$ & $\bH$ & IX \\ 
 $(\bC)$ & $\bC$ & $(0,1)$ & $\bC$  & $\bC$ & $\bC \times \bC$ & VII \\ 
 $\bC[\epsilon]$ & $\bC$ & $(1,1)$ & $\bC$  & $\bC \times \bC$ & $M_2(\bC)$ & IV \\
  \bottomrule
\end{tabular}
\end{center} 
\end{table}

\begin{proof}
    It is clear that all the listed $C_2$-graded algebras $D$ are $C_2$-graded division algebras, and the values of $\widehat{Z}(D)$, $j(D)$, and $D_0$ are direct to compute. It is also clear that no pair are isomorphic as $C_2$-graded algebras, by considering the invariant $j(D)$ and the $\bR$-dimension. Conversely, given any $C_2$-graded division algebra $D$ over $\bR$, this defines an element $[D] \in \BrW(\bR) \sqcup \BrW(\bC)$ as above. These have orders $|\BrW(\bR)| = 8$ and $|\BrW(\bC)| = 2$, and therefore $[D] = [D']$ for some $D'$ on our list, because they all define distinct classes by the injectivity of Corollary \ref{cor:divalgrep} and therefore must exhaust $\BrW(\bR) \sqcup \BrW(\bC)$. But then $D \cong D'$, again by the injectivity of Corollary \ref{cor:divalgrep}.
    
    To understand the group structure of $\BrW(\bR)$ in terms of these representatives, note that there is a $C_2$-graded isomorphism
    \[
    \bR[\epsilon] \, \hat{\otimes} \, \bR[\epsilon] \xrightarrow{\sim} \bC_{\hspace{-0.1em}\rtimes \!}[\epsilon],
    \]
    where $1 \, \hat{\otimes} \, 1$, $1 \, \hat{\otimes} \, \epsilon$, $\epsilon \, \hat{\otimes} \, 1$ and $\epsilon \, \hat{\otimes}\, \epsilon$ map to $1$, $\epsilon$, $i \epsilon$ and $i$ respectively. The graded opposite algebra (the inverse in the Brauer-Wall group) of $\bC_{\hspace{-0.1em}\rtimes \!}[\epsilon]$ is $\bC_{\hspace{-0.1em}\rtimes \!}[\delta]$, which follows from the definition, and therefore $\bC_{\hspace{-0.1em}\rtimes \!}[\epsilon]$ is not self-inverse, hence has order neither $1$ or $2$, so $\bR[\epsilon]$ must have order $8$. This implies that $4 \cdot [\bR[\epsilon]] = [(\bH)]$, and the description of the remaining rows as powers of $\bR[\epsilon]$ follows from noting that $\bR[\delta]$ is the graded opposite algebra of $\bR[\epsilon]$ (so $[\bR[\delta]]$ is the inverse of $[\bR[\epsilon]]$), and that
    \[
        \bH[\epsilon] \cong (\bH) \, \hat{\otimes} \, \bR[\epsilon], \qquad \bH[\delta] \cong (\bH) \, \hat{\otimes} \, \bR[\delta].
    \]

    It remains to compute $D_{\text{ss}}$ and $D_i$. For the ungraded rows, it is clear that $(\bR)_{i} \cong \bC$ and
    \[
    (\bC)_i \cong \bC[x] / (x^2 + 1) \cong \bC \times \bC,
    \]
    and for $(\bH)$ there is an $\bR$-algebra isomorphism
\[
(\bH)_{i} \xrightarrow{\sim} M_2(\bC),
\]
where
\[
\mathbf{i} \mapsto \begin{pmatrix}
    0 & -1 \\ 1 & 0
\end{pmatrix}, \quad \mathbf{j} \mapsto \begin{pmatrix}
    0 & i \\ i & 0
\end{pmatrix}, \quad \mathbf{k} \mapsto \begin{pmatrix}
    -i & 0 \\ 0 & i
\end{pmatrix}, \quad \underline{i} \mapsto \begin{pmatrix}
    i & 0 \\ 0 & i
\end{pmatrix}.
\]

For rows $2, 6$ and $10$, suppose that $A \in \{\bR, \bC, \bH\}$. Then as $Z(A)$-algebras,
\[
A[\epsilon] \xrightarrow{\sim} A \times A, \qquad x + y \epsilon \mapsto (x+y, x-y).
\]
Furthermore,
\[
A[\epsilon]_i = A[\epsilon] \oplus A[\epsilon] \cdot \underline{i} \xrightarrow{\sim} M_2(A),
\]
where
\[
(x + y \epsilon, (u + v \epsilon) \cdot \underline{i}) \mapsto \begin{pmatrix}
    x + y & -u - v \\ u - v & x - y
\end{pmatrix}.
\]

All that remains are rows $3$, $4$, $7$ and $8$. For $\bC_{\hspace{-0.1em}\rtimes \!}[\epsilon]$, there is an isomorphism
\[
\bC_{\hspace{-0.1em}\rtimes \!}[\epsilon] \xrightarrow{\sim} M_2(\bR), \qquad (a+ib + (c + id) \epsilon) \mapsto \begin{pmatrix}
    x + y & -u - v \\ u - v & x - y
\end{pmatrix},
\]
which comes from noting that $\bC_{\hspace{-0.1em}\rtimes \!}[\epsilon]$ is nothing but $\bC_{\star \hspace{-0.1em}}[C_2]$ for the trivial real structure $1 \leq C_2$: the isomorphism is the action map $\bC_{\star \hspace{-0.1em}}[C_2] \xrightarrow{\sim} \End_{\bR}(\bC)$ on the trivial module $\bC$, and taking the basis $\{1, i\}$. To see that $\bC_{\hspace{-0.1em}\rtimes \!}[\epsilon]_i \cong M_2(\bR) \times M_2(\bR)$, we first note that $\bC_{\hspace{-0.1em}\rtimes \!}[\epsilon]_i$ is a semisimple $\bR$-algebra, which follows by \cite[Def.\ IV.3.5, Thm.\ IV.3.6]{LAM} from the fact that it is $C_2$-graded simple (Lemma \ref{lem:Aalphaproperties}). From Lemma \ref{lem:Aalphaproperties} we also have that
\[
Z(\bC_{\hspace{-0.1em}\rtimes \!}[\epsilon]_i) = \bR \oplus \bR \cdot i \cdot \underline{i} \cong \bR[t] / (t^2 - 1) \cong \bR \times \bR
\]
so $\bC_{\hspace{-0.1em}\rtimes \!}[\epsilon]_i \cong A \times A'$ for some central simple $\bR$-algebras $A, A'$. The composition of the inclusion $M_2(\bR) \hookrightarrow \bC_{\hspace{-0.1em}\rtimes \!}[\epsilon]_i$ followed by the projection to either factor is injective (by the simplicity of $M_2(\bR)$), hence $A \cong M_2(\bR) \cong A'$ by dimension considerations.

For $\bC_{\hspace{-0.1em}\rtimes \!}[\delta]$, there is a non-zero $\bR$-algebra homomorphism
\[
\bH \xrightarrow{\sim} \bC_{\hspace{-0.1em}\rtimes \!}[\delta], \qquad \mathbf{i} \mapsto i, \: \mathbf{j} \mapsto \delta, \: \mathbf{k} \mapsto i \delta,
\]
which is injective by the simplicity of $\bH$, hence an isomorphism. Similarly to above, $\bC_{\hspace{-0.1em}\rtimes \!}[\delta]_i$ is a semisimple $\bR$-algebra with centre $\bR \times \bR$ and an $\bR$-subalgebra $\bH \hookrightarrow \bC_{\hspace{-0.1em}\rtimes \!}[\delta]_i$ hence $\bC_{\hspace{-0.1em}\rtimes \!}[\delta]_i \cong \bH \times \bH$.

For $\bR[\delta]$, it is clear that $\bR[\delta] \cong \bC$ as $\bR$-algebras, and the isomorphism $\bR[\delta]_i \cong \bH$, follows because $\bR[\delta]_i$ is nothing but $\bC_{\hspace{-0.1em}\rtimes \!}[\delta]$.

Finally, for $\bH[\delta]$, $\bH[\delta] \cong (\bH)_i$, where $\delta$ corresponds to $\underline{i}$, and therefore $\bH[\delta] \cong M_2(\bC)$ from the case of $(\bH)$ above. To see that $\bH[\delta]_i \cong M_4(\bR)$, note that the centraliser of $\bH$ in $\bH[\delta]_i$ is 
\[
\bR[\delta]_i = \bR[\delta] \oplus \bR[\delta] \cdot \underline{i} \leq \bH[\delta]_i.
\]
We have that $\bR[\delta]_i \cong \bH$ from the case of $\bR[\delta]$ above, and therefore there is an $\bR$-algebra homomorphism
\[
\bH \otimes_{\bR} \bH \rightarrow \bH[\delta]_i,
\]
Because $\bH$ has order two in $\Br(\bR)$, $\bH \otimes_{\bR} \bH \cong M_4(\bR)$, and therefore the $\bR$-algebra homomorphism
\[
M_4(\bR) \xrightarrow{\sim} \bH \otimes_{\bR} \bH \rightarrow \bH[\delta]_i
\]
is an isomorphism, as $M_4(\bR)$ is simple and $\bH[\delta]_i$ has the same $\bR$-dimension.
\end{proof}

We can use this to give a short proof of Dyson's Tenfold Way.

\begin{thm}[Dyson's Tenfold Way]\label{thm:dyson10}
Let $W \in \Irr_{\bC}^{\hspace{0.1em} \star}(G)$. Then the pair $(D(W),\fD(W))$ is as in one of the rows of Table \ref{tab:graddivoverR}, for $D \coloneqq D(W)$. In particular, $W$ cannot have D-type II.V or IX.V, and $\fD(W)$ uniquely determines and is determined by the isomorphism type of $D(W)$ as a $C_2$-graded $\bR$-algebra.
\end{thm}

\begin{proof}
    We have that as $\bR$-algebras, $D(W)_0 = \End_{\bC_{\star \hspace{-0.1em}}[G]}(W)$, $D(W)_{\text{ss}} = \End_{\bR[G]}(W)$, and $D(W)_i \cong \End_{\bR[H]}(W)$ by Corollary \ref{cor:Dalphaisom}, and therefore the triples 
    \[
    (D(W)_0, D(W)_{\text{ss}}, D(W)_i) \:\: \text{ and } \: \: (\End_{\bC_{\star \hspace{-0.1em}}[G]}(W), \End_{\bR[G]}(U), \End_{\bR[H]}(X))
    \]
    have Brauer-equivalent simple factors in each component. If $W$ has D-type II.V, this triple would be $(\bR,\bC,\bR)$, using Table \ref{tab:12types}, which occurs for no line of Table \ref{tab:graddivoverR}, and therefore this D-type cannot occur. Similarly, if $W$ has D-type IX.V, this tripe would be $(\bH, \bC, \bH)$, which does not appear in Table \ref{tab:graddivoverR}.

    For the remaining 10 D-types, we can similarly read off the isomorphism class of $D(W)$ as a $C_2$-graded division algebra. For example, if $W$ has D-type II, then by Table \ref{tab:12types} the triple is $(\bH, \bC, \bR)$, which shows that $D(W) \cong \bH[\delta]$ as $C_2$-graded algebras. This determines $D(W)$ in all D-types apart from when $W$ has D-type IV or VII, which both have triple $(\bC, \bC, \bC)$, and therefore there are two options for $D(W)$ as either $(\bC)$ or $\bC[\epsilon]$. These can be distinguished by noting that $D(W)$ has $D(W)_1 = 0$ if and only if $W \not\cong \overline{W}$ (Proposition \ref{prop:basicpropDW}), which occurs in case VII and not in case IV (using Table \ref{tab:24types5}), and so $D(W) \cong \bC[\epsilon]$ in case IV and $D(W) \cong (\bC)$ in D-type VII.
\end{proof}

\begin{remark}
    All of the ten D-types of Table 4.1 are realised by irreducible antilinear representations of finite groups \cite[Thm.\ 3.6]{RUMTAY}, and therefore any $C_2$-graded division algebra over $\bR$ is realised as $D(W)$ for some irreducible antilinear representation $W$ of a finite group by Theorem \ref{thm:dyson10}.
\end{remark}

We can use Theorem \ref{thm:dyson10} to alternatively characterise the 10 D-types in terms of $A(W)$.

\begin{cor}\label{cor:DysonforAW}
If $W \in \Irr_{\bC}^{\hspace{0.1em} \star}(G)$, then $\widehat{Z}(A(W)) = \widehat{Z}(D(W))$. When $\widehat{Z}(D(W)) = \bR$, 
\[
        [A(W)] \cdot [D(W)] = [\bR[\epsilon]]
\]
in $\BrW(\bR)$. When $\widehat{Z}(D(W)) = \bC$, $[A(W)] = [D(W)]$ in $\BrW(\bC)$.
\end{cor}

\begin{proof}
    By Proposition \ref{prop:relateAWDW}, $\widehat{Z}(A(W)) = \bR$ if and only if $\widehat{Z}(D(W)) = \bR$, in which case
    \[
        [A(W)] \cdot [D(W)] = [\bR[\epsilon]]
    \]
    in $\BrW(\bR)$. That $\widehat{Z}(D(W)) = \bR$ occurs in all D-types but IV and VII, and so it only remains to understand $A(W)$ in D-types IV and VII, for which $\widehat{Z}(A(W)) \neq \bR$.

    When $W$ has D-type IV, $X \cong {}^{\omega} X$ and $U \not\cong U \otimes \pi$ from Table \ref{tab:24types5}, and thus $A(W)_0$ is simple with $Z(A(W)_0) \cong \bC$ and $Z(A(W)) \cong \bC \times \bC$ by Proposition \ref{prop:descofSblock} and Lemma \ref{lem:basicpropertiesAW}. In particular, $\widehat{Z}(A(W)) \leq Z(A(W)_0) \cong \bC$ by Lemma \ref{lem:centreofC2graddiv}, and so $\widehat{Z}(A(W)) \cong \bC$. Because $Z(A(W)) > \widehat{Z}(A(W))$, $A(W)$ has \emph{odd type} (see \cite[Def.\ IV.3.5]{LAM}), and thus $[A(W)] = [\bC[\epsilon]]$ in $\BrW(\bC)$, as $[\bC[\epsilon]]$ is the unique central $C_2$-graded division algebra over $\bC$ of odd type, and being of odd type depends only on the class in $\BrW(\bC)$ \cite[Thm.\ IV.4.3]{LAM}.

    Similarly, when $W$ has D-type VII, $X \not\cong {}^{\omega} X$ and $U \cong U \otimes \pi$ from Table \ref{tab:24types5}, so $A(W)$ is simple with $Z(A(W)) = \bC$ by Proposition \ref{prop:descofSblock} and Lemma \ref{lem:basicpropertiesAW}. Now $\widehat{Z}(A(W)) \leq Z(A(W)) \cong \bC$ by Lemma \ref{lem:centreofC2graddiv}, so $\widehat{Z}(A(W)) \cong \bC$, but now $Z(A(W)) = \widehat{Z}(A(W))$, so $A(W)$ has \emph{even type}, and thus $[A(W)] = [(\bC)]$ in $\BrW(\bC)$.
\end{proof}

\begin{remark}
    The paper \cite{GeiMoo} also computes the class of $A(W)$ and $D(W)$ in the appropriate Brauer-Wall group for each of the ten D-types. What we call $A(W)$ is denote $\sA$ there, and their algebra $\sW$ is our $\widetilde{D}(W)$, which by Proposition \ref{prop:relateAWDW} is isomorphic to $D(W)$.
    
    Whilst it allows us to compute the classes of $A(W)$ and $D(W)$, for us the main interest in Theorem \ref{thm:dyson10} is the fact that it gives a very simple proof of Dyson's Tenfold Way using only the classification of $C_2$-graded division algebras over $\bR$ and Corollary \ref{cor:Dalphaisom}, which holds for general types.
    
    As was observed in \cite[\S 5]{GeiMoo}, one can define a ``double'' Frobenius-Schur indicator,
    \[
    \sF'(W) \coloneqq \frac{\sF(W) + i \tilde{\sF}(W)}{|\sF(W) + i \tilde{\sF}(W)|} \in \{e^{\frac{2 \pi i n}{8}}\}_{n \in \bZ / 8 \bZ},
    \]
    which has $\sF'(W) = 0$ when $W$ has D-type IV or VII, and otherwise determines the D-type of $W$. Now,
    \[
        \sF'(W) = e^{\frac{2\pi i n}{8}}, \: \text{ where $n \in \bZ / 8 \bZ$ is defined by } [D(W)] = n \cdot [\bR[\epsilon]]
    \]
    by Theorem \ref{thm:dyson10} and therefore there is an identification of abelian groups of $\{e^{\frac{2 \pi i n}{8}}\}_{n \in \bZ / 8 \bZ}$ with $\BrW(\bR)$. Considering instead $A(W)$, we obtain a second identification of these groups from the relation
    \[
    \sF'(W) = e^{\frac{2\pi i n}{8}}, \: \text{ where $n \in \bZ / 8 \bZ$ is defined by } [A(W)] = (1-n) \cdot [\bR[\epsilon]]
    \]
    by Corollary \ref{cor:DysonforAW}. The two identifications are related by a reflection in the line $\langle 1 + i \rangle_{\bR}$.
\end{remark}

\section{Comparison between Antilinear Representations and Representations over $\bR$}\label{sect:antilinearvsrealreps}

In this section we review some well-known properties of real-representations and observe to what extent these generalise to the antilinear setting. Throughout $H$ is a finite group.

\subsection{Dimension of complex representations of C-type $\bH$}

\begin{prop}\label{prop:evendim}
    Suppose that $V \in \Irr_{\bC}(H)$ has C-type $\bH$. Then $\dim_{\bC}(V)$ is even.
\end{prop}

We note that, unlike the other results in this section, this doesn't actually require that $H$ is finite.

\begin{proof}
    If $V$ corresponds to $W \in \Irr_{\bR}(H)$, then $\dim_{\bR}(W) = 2 \dim_{\bC}(V)$, and $4 \mid \dim_{\bR}(W)$ because $W$ is a module over $\bH$, hence $\dim_{\bC}(V)$ is even.
\end{proof}

This is not true for antilinear representations, which are now complex vector spaces.

\begin{eg}\label{eg:dim1typeH}
    Let $C_2$ have Real structure $C_2 \leq C_4$, and let $V$ be the non-trivial $1$-dimensional representation of $C_2$. Then, writing $C_2 = \langle x \rangle \leq \langle y \rangle = C_4$ with $y^2 = x$, we many compute
    \[
        \sF(V) = \frac{1}{2} \left( \chi_V(y^2) + \chi_V((y^3)^2) \right) = \frac{1}{2} \left( \chi_V(x) + \chi(x) \right) = -1
    \]
    hence $V$ has S-type $\bH$, but dimension $1$.
\end{eg}

\subsection{Abelian groups have no complex representations of C-type $\bH$}

\begin{cor}
Suppose that $H$ is abelian. Then any $V \in \Irr_{\bC}(H)$ has C-type $\bR$ or $\bC$.
\end{cor}

\begin{proof}
    This follows from Proposition~\ref{prop:evendim} or can be seen directly noting that $\bR[H]$ is commutative, so there can be no copy of $\bH$ in its Artin-Wedderburn decomposition.
\end{proof}

This is not true for antilinear representations: from Example \ref{eg:dim1typeH}, $C_2 \leq C_4$ provides an example of Real structure $H \leq G$ where both $H$ and $G$ are abelian, but there is a representation of S-type $\bH$.

\subsection{Burnside's Theorem}

As a consequence of Proposition \ref{prop:evendim} we have the following.

\begin{cor}\label{cor:weakburnside}
Suppose that $|H|$ is odd. Then any $V \in \Irr_{\bC}(H)$ has C-type $\bR$ or $\bC$.
\end{cor}

\begin{proof}
    If there was $V \in \Irr_{\bC}(H)$ of C-type $\bH$, then $V$ has even dimension by Proposition \ref{prop:evendim}, which is impossible as $\dim_{\bC}(V) \mid |H|$ because $V$ is absolutely irreducible.
\end{proof}

Burnside proved the following stronger statement.

\begin{prop}\label{prop:burnside}
    Suppose that $|H|$ is odd. Then any $V \in \Irr_{\bC}(H)$ which is non-trivial has C-type $\bC$.
\end{prop}

\begin{proof}
    As $H$ has odd order, $h \mapsto h^2$ is injective, hence bijective on $H$, so
    \[
        \sF(V) = \sum_{h \in H} \chi(h^2) = \sum_{h \in H}\chi(h) = \dim_{\bC} V^H
    \]
    which is $0$ apart from the case when $V$ is the trivial representation. Alternatively, one can use \cite[Cor.\ 5.19]{RUMTAY} for the Real structure $H \leq H \times C_2$: the odd order condition means that there is only one odd element which squares to $1$, and hence, given that there are no irreducible characters of C-type $\bH$ by Proposition \ref{prop:evendim}, there also are no non-trivial characters of C-type $\bR$.
\end{proof}

The following example shows that Proposition \ref{prop:burnside} is not true for antilinear representations.

\begin{eg}
Consider $C_3 \leq D_6$, and let $V$ be any non-trivial $1$-dimensional representation of $C_3$ over $\bC$. Writing $C_3 = \langle x \rangle \leq \langle x, y \mid x^3, y^2, yxy^{-1} = x^{-1} \rangle$, we can compute that
\[
    \sF(V) = \frac{1}{3} \left( \chi_V(y^2) + \chi_V((xy)^2) + \chi_V((x^2y)^2) \right)
    = \frac{1}{3} \left( \chi_V(1) + \chi_V(1) + \chi_V(1) \right) = 1,
\]
and so the S-type of any $V \in \Irr_{\bC}(C_3)$ for $C_3 \leq D_6$ is $\bR$.
\end{eg}

\subsection{Order of a group with a quaternionic representation}

One can go even further than Corollary \ref{cor:weakburnside}, which says that if there is some $V \in \Irr_{\bC}(H)$ of C-type $\bH$, then $|H|$ is even.

\begin{prop}\label{prop:strongweakburnside}
    Suppose that there is some $V \in \Irr_{\bC}(H)$ of C-type $\bH$. Then $4 \mid |H|$.
\end{prop}

\begin{proof}
    If instead $|H| = 4k+2$, then any Sylow $2$-subgroup of $H$ is cyclic of order $2$, and the claim follows from \cite[Thm.\ 10.9]{ISA}, which states that if $H$ has a Sylow $2$-subgroup which is elementary abelian then $H$ has no irreducible characters of C-type $\bH$. 
\end{proof}

This Proposition \ref{prop:strongweakburnside} is not true for antilinear representations, as shown by Example \ref{eg:dim1typeH}.

\subsection{Antilinear representations of groups of odd order}

We can however use Proposition \ref{prop:strongweakburnside} to show that Corollary \ref{cor:weakburnside} extends to antilinear representations.

\begin{prop}\label{prop:typesforoddorder}
    Suppose that $H$ has odd order. Then every non-trivial $V \in \Irr_{\bC}(H)$ has D-type IV, V or VII. In particular, $V$ has S-type $\bR$ or $\bC$.
\end{prop}

Note that the proof uses the strong form Proposition \ref{prop:strongweakburnside} of Corollary \ref{cor:weakburnside} from the classical case, and the tenfold way.

\begin{proof}
Because the C-type of $V$ is $\bC$, the D-type of $V$ is one of IV, V, VI or VII. In D-type VI, the $U \in \Irr_{\bR}(G)$ corresponding to $V$ has C-type $\bH$, which can't happen by Proposition \ref{prop:strongweakburnside} because the order of the Real structure $G$ is not a multiple of $4$.
\end{proof}

\section{Realisability over Cyclotomic Fields}\label{sect:realisablecyclo}

We continue in this section with the assumption that $H$ is finite. Let $e \coloneqq \exp(H)$, the exponent of $H$, when $\exp(H) > 2$ and $e = 4$ when $\exp(H) \leq 2$, and set $E \coloneqq \bQ(\zeta_e)$ to be the cyclotomic field generated by a primitive $e$th root of unity. The modification when $\exp(H) \leq 2$ is to ensure that $e > 2$, and therefore that the maximal totally real subfield
\[
E_{+} \coloneqq \bR \cap E = \bQ(\zeta_e + \zeta_e^{-1})
\]
has index two in $E$, so that we can consider semilinear $E$-representations for a Real structure $H \leq G$.

The main result of \cite{PAS} is that irreducible $\bR$-representations always have models over $E_+$. Note that when $\exp(H) = 2$, $E_+ = \bQ(i) \cap \bR = \bQ(\zeta_2) \cap \bR = \bQ$, and so our $E_+$ agrees with that of \cite{PAS}.

\begin{thm}\label{thm:classicalcyclotomic}
    Suppose that $W \in \Irr_{\bR}(H)$. Then there is $\tilde{W} \in \Irr_{E_+}(H)$ with $\tilde{W}_{\bR} \cong W$.
\end{thm}

Concretely, this means that any matrix representation of $W$ over $\bR$ can always be conjugated over $\bR$ to matrices which all have values in $E_{+}$. We give the proof for the readers convenience, adapted from that of \cite{PAS} into a form which will be easier for us to generalise to antilinear representations.

\begin{proof}
Let $V \in \Irr_{\bC}(H)$ correspond to $W$, and let $\tilde{V} \in \Irr_{E}(H)$ be a model of $V$ over $E$, so $\tilde{V}_{\bC} \cong V$ and $\chi_{\tilde{V}} = \chi_V$, which always exists because $E$ is a splitting field for $H$ \cite[\S12.3]{SERRE}.

If $V$ has C-type $\bC$ or $\bH$, then $\chi_W = \chi_V + \overline{\chi_V}$. Setting $\tilde{W} \coloneqq \tilde{V}|_{E_+}$, we have that $\chi_{\tilde{W}} = \chi_V + \overline{\chi_V}$, and thus $\tilde{W}_{\bR} \cong W$ by comparing characters.

We may therefore assume that $V$ has C-type $\bR$. There are two options for $\overline{m}_{E_+}^E(\tilde{V}) \in \{1,2\}$. If $\overline{m}_{E_+}^E(\tilde{V}) = 1$, then there is some $\tilde{W} \in \Irr_{E_+}(H)$ with $\tilde{W}_{E} \cong V$, so $\chi_{\tilde{W}} = \chi_V = \chi_W$, hence $\tilde{W}_{\bR} \cong W$ and we're done.

Suppose instead, for a contradiction, that $\overline{m}_{E_+}^E(\tilde{V}) = 2$. The key fact we use is that this means that $2 \mid \langle \chi_W , \chi_U \rangle$ for any $U \in \Rep_{E_+}(H)$ by \cite[Cor.\ 10.3(4)]{TAYSchur} and the fact that $\chi_W = \chi_V = \chi_{\tilde{V}}$. The second input is Serre's induction theorem \cite[Lem.\ 73.19]{CR2}, which says that we can write
    \[
        \chi_W = \sum_{i=1}^r a_i \cdot \Ind_{H_i}^H \psi_i
    \]
    where $r \geq 1$, $a_i \in \bZ$, each $H_i$ is a subgroup of $H$, and $\psi_i$ is a character of $W_i \in \Irr_{\bR}(H_i)$ such that $W_i$ is inflated from a cyclic or dihedral group. Because $W$ is irreducible of C-type $\bR$,
    \[
        1 = \langle \chi_W, \chi_W \rangle = \sum_{i = 1}^r a_i \langle \chi_W, \Ind_{H_i}^H \psi_i \rangle.
    \]
    
    Now each $W_i$ admits a model over $E_+$, as the irreducible $\bR$-representations of cyclic and dihedral groups are defined over $E_+$, and therefore $\Ind_{H_i}^H \psi_i$ is also the character of an object of $\Rep_{E_+}(H)$. This means that each $\langle \chi_W, \Ind_{H_i}^H \psi_i \rangle$ is even, and so $1$ is a sum of even numbers, a contradiction.
\end{proof}

The goal of this section is to prove the analogous statement for antilinear representations.

\begin{thm}\label{thm:antilinearcyclotomicfinal}
    Suppose $W \in \Irr_{\bC}^{\hspace{0.1em} \star}(G)$. Then there is $\tilde{W} \in \Irr_{E}^{\hspace{0.1em} \star}(H)$ with $\tilde{W}_{\bC} \cong W$.
\end{thm}

Concretely, this is equivalent to the claim that any antilinear matrix representation of $G$ with values in $\bC$ can always be (antilinear) conjugated over $\bC$ to matrices which all have values in $E$. In the case of split extensions this recovers Theorem \ref{thm:classicalcyclotomic}.

Let us note two cases when Theorem \ref{thm:antilinearcyclotomicfinal} is immediate. The first is when $W$ is $1$-dimensional.

\begin{prop}
Suppose that $V \in \Irr_E(H)$ is $1$-dimensional, and extends to a antilinear representation of $G$ over $\bC$. Then $V$ extends to a antilinear representation of $G$ over $E$.
\end{prop}

\begin{proof}
    We use the matrix form of semilinear representations from Section \ref{sect:matrixform} of Chapter \ref{chapter:algebras}. In this framework, a $1$-dimensional antilinear representation over $\bC$ can be described by a group homomorphism $\rho \colon G \rightarrow \GL_1^{\star}(\bC)$, where $\GL_1^{\star}(\bC) = \GL_1^1(\bC) \sqcup \GL_1^{-1}(\bC)$, where elements are written as $(a)_1$ and $(a)_{-1}$ respectively for $a \in \bC^{\times}$. Multiplication is given by
    \[
    (a)_{1} \cdot (b)_{\epsilon} = (ab)_{\epsilon} \: \text{ and } \: (a)_{-1} \cdot (b)_{\epsilon} = (a \overline{b})_{-\epsilon}.
    \]
    
    If $\rho(g) = (r_g)_{\epsilon(g)}$, then the equivalent matrix representations are $\rho_x$ for some $x \in \bC^{\times}$, where
    \[
        \rho_x(g) = \twopartdef{(r_g)_1}{\epsilon = 1}{(\overline{x} x^{-1} r_g)_{-1}}{\epsilon = -1}.
    \]

    The map $\bC^{\times} \rightarrow \bC^{\times}$, $x \mapsto \overline{x} x^{-1}$ has image $S^1$, and so the equivalent representations are parametrised by $y \in S^1$, where
    \[
    \rho^y(g) = \twopartdef{(r_g)_1}{\epsilon(g) = 1}{(y r_g)_{-1}}{\epsilon(g) = -1}.
    \]

    Let $\chi \colon H \rightarrow E^{\times}$ be the character of $V$, and suppose that $\rho \colon G \rightarrow \GL_1^{\star}(\bC)$ extends $\chi$, meaning that $r_g = \chi(g)$ for all $g \in H$. To show that $\chi$ extends to an antilinear representation over $E$ is to ask that one can find $y \in S^1$ with $\rho^y$ valued in $\GL_1^{\star}(E)$, meaning that all $y r_g \in E^{\times}$.

    For any $g_1, g_2 \in G \setminus H$, $g_1 g_2 \in H$, so $(r_{g_1} \overline{r_{g_2}})_1 = \rho(g_1g_2) = (r_{g_1 g_2})_1$. We know that $r_{g_1 g_2}$ is an $e$th root of $1$, and therefore (as $e$th roots are closed under multiplication and complex conjugation) $r_{g_1}$ is an $e$th root of $1$ if and only if $r_{g_2}$ is an $e$th root of $1$.
    
    In particular, we may fix any $\omega \in G \setminus H$, and set $y = r_\omega^{-1}$. Then $r^y(\omega) = (1)_{-1}$, and $\rho^y$ is valued in $\GL_1^{\star}(E)$. To see that this is well-defined, meaning that $y \in S^1$, note that because $g^2 \in H$, $|r_g| = r_g \overline{r_g}$ is an $e$th root of $1$. We also have that $r_g \overline{r_g} \in \bR_{> 0}$, and therefore $|r_g| = 1$.
\end{proof}

We second case when the conclusion of Theorem \ref{thm:antilinearcyclotomicfinal} is immediate is when $W$ has S-type $\bC$ or $\bH$.

\begin{lemma}\label{lem:onlycareabouttypeR}
    Suppose $W \in \Irr_{\bC}^{\hspace{0.1em} \star}(G)$ has S-type $\bC$ or $\bH$. Then there is $\tilde{W} \in \Irr_{E}^{\hspace{0.1em} \star}(H)$ with $\tilde{W}_{\bC} \cong W$. 
\end{lemma}

\begin{proof}
    If $V \in \Irr_{\bC}(H)$ corresponds to $W$ and has S-type $\bC$ or $\bH$, then $\chi_W = \chi_V + {}^{\omega}\overline{\chi_V}$. Setting $\tilde{W} \coloneqq \Ind_H^G(\tilde{V})$, where $\tilde{V}$ is a model of $V$ over $E$, we have that 
    \[
    \chi_{\tilde{W}} = \chi_{\tilde{V}} + {}^{\omega}\overline{\chi_{\tilde{V}}} = \chi_V + {}^{\omega}\overline{\chi_V}
    \]
    hence $\tilde{W}_{\bC} \cong W$ as they have the same character.
\end{proof}

Therefore, to prove Theorem \ref{thm:antilinearcyclotomicfinal}, one must only consider the case when $W$ has D-type I, V, or IX. To understand the potential obstructions in these cases, let $V \in \Irr_{\bC}(G)$ correspond to $W$, and let $\tilde{V}$ be a model of $V$ over $E$. If $\tilde{W} \in \Irr_{E}^{\hspace{0.1em} \star}(G)$ corresponds to $\tilde{V}$, then we either have that $\tilde{W}_{\bC} \cong W$ (when $\tilde{W}$ has S-type $\bR$) or $\tilde{W}_{\bC} \cong W \oplus W$ (when $\tilde{W}$ has S-type $\bH$). To prove Theorem \ref{thm:antilinearcyclotomicfinal} is to prove that this second case cannot occur.

The classification into 24 ways alone is not enough to deduce this. For example, a priori, the pair $(V,W)$ could have D-type I, but $(\tilde{V}, \tilde{W})$ have type (9). Similarly, $(V,W)$ could have D-type V, and $(\tilde{V}, \tilde{W})$ have type (15), or $(V,W)$ could have D-type IX, whilst $(\tilde{V}, \tilde{W})$ have type (19).

Instead, in order to prove Theorem \ref{thm:antilinearcyclotomicfinal}, we will follow the general strategy of the classical proofs for cyclotomic realisability over $\bC$ and $\bR$, and deduce this as a consequence of a semilinear induction theorem for antilinear representations.

\subsection{Fundamental Representations}\label{sect:fundreps}

There are two natural operations one can consider on antilinear representations. The first is inflation, and the second is induction from a subgroup.

In this section we consider families of Real groups which (in a precise sense) generate all semilinear representations of all Real groups using these two operations. Throughout this section, $\bL / \bK$ is a characteristic $0$ degree two Galois extension, which we will later take to be $\bC / \bR$.

It turns out that it is most natural to work with the category of $C_2$-graded groups, which we consider as pairs $(G, \mu_G \colon G \rightarrow C_2)$, where a morphism $(G, \mu_G) \rightarrow (H,\mu_H)$ is a group homomorphism $f \colon G \rightarrow H$ such that $\mu_H \circ f = \mu_G$. The category of Real groups is recovered as the full subcategory which has objects those $(G,\mu_G)$ such that $\mu_G \colon G \rightarrow C_2$ is surjective.

A $C_2$-graded subgroup (resp.\ quotient) is simply a morphism of $C_2$-graded groups which is injective (resp.\ surjective). For a subgroup $H$ of a $C_2$-graded group $G$, $H$ has a canonical $C_2$-graded group structure, which is the unique structure making the inclusion a morphism of $C_2$-graded groups. For a quotient $f \colon G \twoheadrightarrow Q$ of $C_2$-graded group, $Q$ has at most one structure as a $C_2$-graded group making $f \colon G \twoheadrightarrow Q$ a morphism of $C_2$-graded groups, which exists if and only if $\ker(f) \subset \ker(\mu_G)$.

For a $C_2$-graded group $G$, we write $\bL_{\star \hspace{-0.1em}}[G]$ for the skew group algebra of $G$, which has our standard meaning when $\mu_G \colon C \rightarrow C_2 \equiv \Gal(\bL / \bK)$ is surjective, but which we also use to denote the ordinary group algebra $\bL[G]$ when $\mu_G$ is trivial. Similarly, we extend the notation $\Rep_{\bbL}^{\hspace{0.1em} \star}(G)$ and $\Irr_{\bbL}^{\hspace{0.1em} \star}(G)$ to the case where $\mu_G$ is potentially trivial. With this definition, $\bL_{\star \hspace{-0.1em}}[-]$ can be viewed as a functor from $C_2$-graded groups to $\bK$-algebras. In particular, for a $C_2$-graded subgroup $H \hookrightarrow G$ we can define $C_2$-graded induction by
\[
\Ind_S^G(-) \coloneqq \bL_{\star \hspace{-0.1em}}[G] \otimes_{\bbL_{\star \hspace{-0.1em}}[H]} -,
\]
and for a $C_2$-graded quotient $f \colon G \twoheadrightarrow Q$, we can define $C_2$-graded inflation $\Inf_Q^G(-)$ as restriction along the natural map $\bL_{\star \hspace{-0.1em}}[G] \rightarrow \bL_{\star \hspace{-0.1em}}[Q]$.

\begin{defn}
    A $C_2$-graded subquotient of $(G, \mu_G)$ is a $C_2$-graded quotient of a $C_2$-graded subgroup of $G$.
\end{defn}

The following shows that the $C_2$-graded subquotient relation is transitive and establishes a property of representations inflated and induced from $C_2$-graded subquotients, which we will use later.

\begin{prop}\label{prop:subquotofsubqout}
    Suppose that we have a diagram of $C_2$-graded groups:
    % https://q.uiver.app/#q=WzAsNSxbMCwwLCJHIl0sWzAsMSwiSCJdLFsxLDEsIlEiXSxbMSwyLCJTIl0sWzIsMiwiUCJdLFsxLDIsImYiLDAseyJzdHlsZSI6eyJoZWFkIjp7Im5hbWUiOiJlcGkifX19XSxbMSwwLCIiLDIseyJzdHlsZSI6eyJ0YWlsIjp7Im5hbWUiOiJob29rIiwic2lkZSI6InRvcCJ9fX1dLFszLDIsIiIsMix7InN0eWxlIjp7InRhaWwiOnsibmFtZSI6Imhvb2siLCJzaWRlIjoidG9wIn19fV0sWzMsNCwiZyIsMCx7InN0eWxlIjp7ImhlYWQiOnsibmFtZSI6ImVwaSJ9fX1dXQ==
\[\begin{tikzcd}
	G \\
	H & Q \\
	& S & P
	\arrow[hook, from=2-1, to=1-1]
	\arrow["f", two heads, from=2-1, to=2-2]
	\arrow[hook, from=3-2, to=2-2]
	\arrow[two heads, from=3-2, to=3-3]
\end{tikzcd}\]
Then viewing $f^{-1}(S)$ as a $C_2$-graded group using the inclusion $f^{-1}(S) \hookrightarrow G$, the map $f^{-1}(S) \twoheadrightarrow P$ is a $C_2$-graded quotient, and there is a natural isomorphism
\[
\Inf_Q^H(\Ind_S^Q(\Inf_P^S(V))) \cong \Ind^H_{f^{-1}(S)}(\Inf_P^{f^{-1}(S)}(V))
\]
for any $V \in \Rep_{\bbL}^{\hspace{0.1em} \star}(P)$.
\end{prop}

\begin{proof}
    The first claim follows from the commutativity of the diagram
% https://q.uiver.app/#q=WzAsNyxbMCwxLCJIIl0sWzEsMSwiUSJdLFsxLDIsIlMiXSxbMCwyLCJmXnstMX0oUykiXSxbMiwwLCJDXzIiXSxbMCwwLCJHIl0sWzIsMiwiUCJdLFswLDEsImYiLDAseyJzdHlsZSI6eyJoZWFkIjp7Im5hbWUiOiJlcGkifX19XSxbMiwxLCIiLDIseyJzdHlsZSI6eyJ0YWlsIjp7Im5hbWUiOiJob29rIiwic2lkZSI6InRvcCJ9fX1dLFszLDIsImYiLDAseyJzdHlsZSI6eyJoZWFkIjp7Im5hbWUiOiJlcGkifX19XSxbMywwLCIiLDEseyJzdHlsZSI6eyJ0YWlsIjp7Im5hbWUiOiJob29rIiwic2lkZSI6InRvcCJ9fX1dLFsxLDQsIlxcbXVfUSJdLFsyLDQsIlxcbXVfUyIsMix7ImN1cnZlIjoxfV0sWzAsNCwiXFxtdV9IIiwwLHsiY3VydmUiOi0xfV0sWzUsNCwiXFxtdV9HIiwwLHsiY3VydmUiOi0xfV0sWzAsNSwiIiwwLHsic3R5bGUiOnsidGFpbCI6eyJuYW1lIjoiaG9vayIsInNpZGUiOiJ0b3AifX19XSxbMiw2LCIiLDAseyJzdHlsZSI6eyJoZWFkIjp7Im5hbWUiOiJlcGkifX19XSxbNiw0LCJcXG11X1AiLDIseyJjdXJ2ZSI6MX1dXQ==
\[\begin{tikzcd}
	G && {C_2} \\
	H & Q \\
	{f^{-1}(S)} & S & P
	\arrow["{\mu_G}", curve={height=-6pt}, from=1-1, to=1-3]
	\arrow[hook, from=2-1, to=1-1]
	\arrow["{\mu_H}", curve={height=-6pt}, from=2-1, to=1-3]
	\arrow["f", two heads, from=2-1, to=2-2]
	\arrow["{\mu_Q}", from=2-2, to=1-3]
	\arrow[hook, from=3-1, to=2-1]
	\arrow["f", two heads, from=3-1, to=3-2]
	\arrow["{\mu_S}"', curve={height=6pt}, from=3-2, to=1-3]
	\arrow[hook, from=3-2, to=2-2]
	\arrow[two heads, from=3-2, to=3-3]
	\arrow["{\mu_P}"', curve={height=6pt}, from=3-3, to=1-3]
\end{tikzcd}\]

For the second, given $V \in \Rep_{\bbL}^{\hspace{0.1em} \star}(P)$, we define
\[
\phi \colon \bL_{\star \hspace{-0.1em}}[H] \otimes_{\bbL_{\star \hspace{-0.1em}}[f^{-1}(S)]} \Inf_P^{f^{-1}(S)}(V) \rightarrow \Inf_Q^H(\bL_{\star \hspace{-0.1em}}[Q] \otimes_{\bbL_{\star \hspace{-0.1em}}[S]} \Inf^S_P(V)), \qquad \phi(h \otimes v) \coloneqq f(h) \otimes v,
\]
which is well-defined and surjective. As both sides have the same $\bK$-dimension $\dim_{\bbK}(V) \cdot [H:f^{-1}(S)] = \dim_{\bbK}(V) \cdot [Q:S]$, $\phi$ is an isomorphism.
\end{proof}
\begin{defn}
We write $\Char^{\hspace{0.1em} \star}(G,\bL)$ for the collection of characters of representations of the $C_2$-graded group $G$ over $\bL$.
\end{defn}

\begin{defn}
For a list $\sL$ of isomorphism classes of $C_2$-graded groups, we say that $\sL$ is \emph{generating} if $\sL$ \emph{satisfies condition} $(G)$: for any $C_2$-graded group $G$ and $\chi \in \Char^{\hspace{0.1em} \star}(G, \bL)$ one can write 
\[
\chi = \sum_{i} n_i \cdot \Ind_{G_i}^G(\Inf_{Q_i}^{G_i}(\psi_i)),
\]
where $n_i \in \bZ$, $G_i$ are $C_2$-graded subgroups of $G$, $G_i \twoheadrightarrow Q_i$ are $C_2$-graded quotients, $\psi_i \in \Char^{\hspace{0.1em} \star}(Q_i, \bL)$, and each $Q_i \in \sL$.

We say that such a list $\sL$ is \emph{minimal} if after removing any element of $\sL$, $\sL$ no longer satisfies $(G)$.
\end{defn}

\begin{defn}\label{def:orderC2gradedgroup}
We define the \emph{order} $o(G)$ of a $C_2$-graded group $G$ by 
\[
o(G) \coloneqq 2 |\ker(\mu_G)| + |\mu_G(G)| - 2 \in \bZ_{\geq 1}.
\]
\end{defn}

Concretely, $o(G) = 2|G| - 1$ if $G$ is ungraded, and $o(G) = |G|$ if $G$ is non-trivially graded. With this definition, the order of a $C_2$-graded subgroup (resp.\ $C_2$-graded quotient) of $G$ is always at most the order of $G$, with equality if and only if the inclusion map (resp.\ the quotient map) is an isomorphism. For example, when $G$ is non-trivially graded the proper inclusion $\ker(\mu_G) \hookrightarrow G$ is between groups of order $|G| - 1$ and $|G|$ respectively.

Given a list of isomorphism classes of $C_2$-graded groups $\sL$ and $n \geq 1$, we set $\sL_{\leq n}$ for the subset of $C_2$-graded groups of order $\leq n$, and $\sL_n$ for the subset of $C_2$-graded groups of order exactly $n$.

\begin{prop}\label{prop:minimallistexists}
    A minimal list $\sG$ satisfying \emph{$(G)$} exists, and can be described inductively:
    \[
    \sG_1 = \{(C_1, C_1 \hookrightarrow C_2)\},
    \]
    and for any $n \geq 1$, $\sG_{n+1}$ is defined as the set of $C_2$-graded groups of order $n+1$ for which there is some $\chi \in \Char^{\hspace{0.1em} \star}(G, \bL)$ such that you cannot write
    \[
    \chi = \sum_{i} n_i \cdot \Ind_{G_i}^G(\Inf_{Q_i}^{G_i}(\psi_i)),
    \]
    where $n_i \in \bZ$, $G_i$ are $C_2$-graded subgroups of $G$, $G_i \twoheadrightarrow Q_i$ are $C_2$-graded quotients, $\psi_i \in \Char^{\hspace{0.1em} \star}(Q_i, \bL)$, and each $Q_i \in \sG_{\leq n}$. 
\end{prop}

\begin{proof}
    It is clear that $\sG$, as defined inductively above, satisfies $(G)$. To see that it is minimal, let $G \in \sG_{n+1}$, and suppose for a contradiction that $\sG \setminus \{G\}$ satisfied $(G)$. In particular, for any character $\chi \in \Char^{\hspace{0.1em} \star}(G, \bL)$ one can write
\[
\chi = \sum_{i} n_i \cdot \Ind_{G_i}^G(\Inf_{Q_i}^{G_i}(\psi_i)),
\]
where $n_i \in \bZ$, $G_i$ are $C_2$-graded subgroups of $G$, $G_i \twoheadrightarrow Q_i$ are $C_2$-graded quotients, $\psi_i \in \Char^{\hspace{0.1em} \star}(Q_i, \bL)$, and each $Q_i \in \sG \setminus \{G\}$. Any proper $C_2$-graded subquotient of $G$ has order $\leq n$, and thus each $Q_i \in \sG_{\leq n}$. However this contradictions the fact that $G \in \sG_{n+1}$.
\end{proof}

\begin{prop}\label{prop:minimallistexistsuniquely}
    A minimal list $\sG$ satisfying \emph{$(G)$} exists and is unique. The list $\sG$ can be explicitly described as the set of $C_2$-graded groups $G$ for which there is some $\chi \in \Char^{\hspace{0.1em} \star}(G, \bL)$ such that you cannot write
    \[
    \chi = \sum_{i} n_i \cdot \Ind_{G_i}^G(\Inf_{Q_i}^{G_i}(\psi_i)),
    \]
    for $n_i \in \bZ$, $G_i$ are $C_2$-graded subgroups of $G$, $G_i \twoheadrightarrow Q_i$ are $C_2$-graded quotients, $\psi_i \in \Char^{\hspace{0.1em} \star}(Q_i, \bL)$, and each $Q_i$ is a proper subquotient $G$.

    Furthermore, any list $\sL$ of Real groups satisfying \emph{$(G)$} contains $\sG$.
\end{prop}

\begin{remark}
    The inductive definition of Proposition \ref{prop:minimallistexistsuniquely} is computationally useful if one wants to describe $\sG$ by an inductive search, as one does not need to check condition $(G)$ over all subquotients, but only over those already in your list. The explicit definition of $\sG$ is theoretically useful, as it allows you to determine if certain infinite families are in $\sG$, without a complete knowledge of $\sG$.
\end{remark}

\begin{proof}
    Let $\sH$ denote the list of $C_2$-graded groups defined via the explicit description above. It is clear that for any list $\sL$ of $C_2$-graded groups satisfying $(G)$, $\sH \subset \sL$. We will show that $\sH$ satisfies $(G)$. From this it is immediate that any minimal list $\sL$ satisfying $(G)$ has to be equal to $\sH$, so in particular $\sG = \sH$, and that any list of $C_2$-graded groups satisfying $(G)$ contains $\sG$.

    We prove that $\sH$ satisfies $(G)$ by induction on the order of the $C_2$-graded group $G$ in condition $(G)$. If $o(G) = 1$, then by definition $G \in \sH$, and so $(G)$ is satisfied. Suppose now that $n \geq 1$ and $\sH$ satisfies $(G)$ for all groups of order $\leq n$, and take $G$ a $C_2$-graded group with $o(G) = n + 1$. If $G \in \sH$ then $(G)$ is trivially satisfied, so suppose that $G \not\in \sH$, meaning that for all $\chi \in \Char^{\hspace{0.1em} \star}(G, \bL)$
\[
\chi = \sum_{i} n_i \cdot \Ind_{G_i}^G(\Inf_{Q_i}^{G_i}(\psi_i)),
\]
where $n_i \in \bZ$, $G_i$ are $C_2$-graded subgroups of $G$, $G_i \twoheadrightarrow Q_i$ are $C_2$-graded quotients, $\psi_i \in \Char^{\hspace{0.1em} \star}(Q_i, \bL)$, and each $Q_i$ is a proper subquotient of $G$. For each $i$, because $Q_i$ is a proper subquotient, $o(Q_i) \leq n$, and therefore by our induction hypothesis we can write 
\[
\psi_i = \sum_{j} m_j \Ind_{S_j}^{Q_i}(\Inf_{P_j}^{S_j}(\mu_j)),
\]
where each $P_j \in \sH$. Using Proposition \ref{prop:subquotofsubqout}, if $f \colon H_i \twoheadrightarrow Q_i$ is the quotient map, then $f^{-1}(S_j)$ is a $C_2$-graded subgroup of $G$, $P_j$ is a $C_2$-graded quotient of $f^{-1}(S_j)$, and
\[
\Inf_{Q_i}^{H_i}(\Ind_{S_j}^{Q_i}(\Inf_{P_j}^{S_j}(\mu_j))) \cong \Ind^{H_i}_{f^{-1}(S_j)}(\Inf_{P_j}^{f^{-1}(S_j)}(\mu_j)).
\]
In particular, 
\begin{align*}
\Ind_{G_i}^G(\Inf_{Q_i}^{G_i}(\psi_i)) 
&= \sum_{j} m_j \Ind_{G_i}^G(\Inf_{Q_i}^{G_i}(\Ind_{S_j}^{Q_i}(\Inf_{P_j}^{S_j}(\mu_j)))), \\
&= \sum_{j} m_j \Ind_{G_i}^G(\Ind^{G_i}_{f^{-1}(S_j)}(\Inf_{P_j}^{f^{-1}(S_j)}(\mu_j))), \\
&= \sum_{j} m_j \Ind^{G}_{f^{-1}(S_j)}(\Inf_{P_j}^{f^{-1}(S_j)}(\mu_j)).
\end{align*}
Replacing each $\Ind_{G_i}^G(\Inf_{Q_i}^{G_i}(\psi_i))$ in this way, we can write  
\[
\chi = \sum_{i} a_i \cdot \Ind_{G_i'}^G(\Inf_{Q_i'}^{G_i'}(\psi_i')),
\]
where $a_i \in \bZ$, $G_i'$ are $C_2$-graded subgroups of $G$, $G_i' \twoheadrightarrow Q_i'$ are $C_2$-graded quotients, $\psi_i' \in \Char^{\hspace{0.1em} \star}(Q_i', \bL)$, and each $Q_i' \in \sH_{\leq n}$. In particular, $\sH$ satisfies $(G)$ for $C_2$-graded groups of order $n+1$, and we are done by induction.
\end{proof}

Therefore, we see that for any characteristic $0$ degree two Galois extension $\bL / \bK$, there is a unique minimal list $\sG_{\bL / \bK}$ satisfying $(G)$. We refer to this $\sG_{\bL / \bK}$ as the list of \emph{fundamental representations}. For any given $\bL / \bK$, the main question is whether or not $\sG_{\bL / \bK}$ admits a reasonable description.

\subsection{Dihedral Representations}

The most important groups for our description of $\sG_{\bC / \bR}$ will be dihedral representations, that is, representations of the Real structure $C_n \leq D_{2n}$. Here we consider $C_n = \langle x \rangle \subset D_{2n} = \langle x ,y \mid x^n, y^2, yxy^{-1} = x^{-1} \rangle$, which makes sense whenever $n \geq 1$ (and so for us $D_2 \coloneqq C_2$ and $D_4 \coloneqq C_2 \times C_2$).

Let us first see that the antilinear representations of $C_n \leq D_{2n}$ are always defined over $\bQ(\zeta_n)$.

\begin{lemma}\label{lem:dihedralrealisable}
    Let $n \geq 1$. Then every irreducible representation of $C_n$ over $\bQ(\zeta_n)$ extends to an antilinear representation of the Real structure $C_n \leq D_{2n}$ over $\bQ(\zeta_n)$.
\end{lemma}

\begin{proof}
    Writing $D_{2n} = \langle x, y \mid x^n, y^2, yxy^{-1} = a \rangle$, if $x$ acts on $F$ by multiplication by the $n$th root of unity $\zeta$, then one can let $y$ act on $F$ by complex conjugation.
\end{proof}

Conversely, the conceptual reason why dihedral representations are important for the antilinear representation theory of all Real groups is the following.

\begin{lemma}\label{lem:1dimfftypeRisdihedral}
    Suppose that $H \leq G$ is a Real group with a $1$-dimensional faithful antilinear representation. Then $(H \leq G) \cong (C_n \leq D_{2n})$.
\end{lemma}

In the proof, we fix a basis an view antilinear representations as morphisms of Real groups $G \rightarrow \GL_n^{*}(\bC)$ (as in \cite[\S 2.2]{RUMTAY}), where
\[
\GL_n^*(\bC) = \{A_1 \mid A \in \GL_n(\bC)\} \sqcup \{A_{-1} \mid A \in \GL_n(\bC)\}
\]
with multiplication 
\[
A_{\epsilon}B_{\delta} = (A \cdot {}^{\epsilon}B)_{\epsilon\delta} = \twopartdef{(A B)_{\epsilon\delta}}{\epsilon = 1}{(A \overline{B})_{\epsilon\delta}}{\epsilon = -1.}
\]
\begin{proof}
    Let $\rho \colon G \rightarrow \GL_1^{*}(\bC)$ be a $1$-dimensional faithful antilinear representation of $G$. Then $\rho|_H \colon H \rightarrow \GL_1^1(\bC)$ is faithful, which means that $H$ is cyclic, so $H = \langle x \rangle$ for some $x \in H$. Let $y \in G$ be any odd element, and write $\rho(x) = [a]_1 \in \GL_1(\bC)$ and $\rho(y) = [b]_{-1} \in \GL_1^{-1}(\bC)$. 
    
    If $H$ has order $n$, then $a^n = 1$ in $\bC^{\times}$. Because $yxy^{-1} = x^i$ for some $i \in \bZ / n \bZ$, $yx = x^{i}y$, hence
    \[
    [b a^{-1}]_{-1} = [b \overline{a}]_{-1} = [b]_{-1} [a]_1 = \rho(yx) = \rho(x^{i}y) = [a^i]_1[b]_{-1} = [a^i b]_{-1},
    \]
    and so $a^{i} = a^{-1}$. In particular, $\rho(x^i) = \rho(x^{-1})$, hence by faithfulness $x^i = x^{-1}$ and so $i = -1$.
    
    We also have that $y^2 \in H$, hence $y^2 = x^j$ for some $j \in \bZ / n \bZ$, and
    \[
    [b \overline{b}]_1 = [b]_{-1}[b]_{-1} = \rho(y^2) = \rho(x^j) = [a^j]_1.
    \]
    so $b \overline{b} = a^j$. Now $a^j$ is a root of $1$, and $b \overline{b} \in \bR_{> 0}$, hence $\rho(y^2) = [1]_1$, and thus by faithfulness $y^2 = 1$.
    
    In summary, we have that $G = \langle x, y \rangle$, $G$ has order $2n$, and $G$ satisfies the relations $x^n = 1$, $y^2 = 1$, $yxy^{-1} = x^{-1}$, hence $G \cong D_{2n}$ as Real groups.
\end{proof}

\subsection{A Semilinear Induction Theorem}

The aim of this section is to describe the list $\sG_{\bC / \bR}$ of fundamental representations, in the special case when $\bL / \bK$ is $\bC / \bR$. Because induction and inflation both preserve cyclotomic realisability, this will allow us to prove cyclotomic realisability for all antilinear representations as an immediate consequence, in Section \ref{sect:proofofcyclorealise} below.

\begin{thm}\label{thm:fundamentalrepsexplicit}
    The list $\sG_{\bC / \bR} = \sG_{\bC / \bR}^{\, 1} \sqcup \sG_{\bC / \bR}^{\,-1}$ consists of the ungraded groups
    \[
       \sG_{\bC / \bR}^{\, 1} =  \{C_n \mid n \geq 1, n \neq 2 \!\!\! \mod 4\}
    \]
    and the Real groups
    \[
       \sG_{\bC / \bR}^{\, -1} = \{(C_n \leq D_{2n}) \mid n \geq 1, n \neq 2 \!\!\! \mod 4 \}.
    \]
\end{thm}

In order to prove Theorem \ref{thm:fundamentalrepsexplicit}, we need to choose a class of Real groups which is big enough to generate all antilinear representations of all Real groups under induction, but small enough that all representations are inflated from dihedral representations.

\begin{defn}\label{def:realelementary}
    A \emph{Real-elementary} group is a Real group of the form $C \rtimes P_0 \leq C \rtimes P$, where
    \begin{itemize}
        \item $C$ is cyclic of odd order,
        \item $P_0 \leq P$ is a Real $2$-group,
        \item $P_0$ acts on $C$ trivially, and $P \setminus P_0$ acts on $C$ by $(-)^{-1}$.
    \end{itemize}
\end{defn}

\begin{prop}\label{prop:realelementarygrpsaremonomial}
    Suppose that $C \rtimes P$ is a Real-elementary group. Then every irreducible antilinear representation of $C \rtimes P$ is induced from a $1$-dimensional representation of a subgroup.
\end{prop}

\begin{proof}
    Let $W$ be an irreducible antilinear representation of $C \rtimes P$. If $W$ has S-type $\bC$ or $\bH$, then $W \cong \Ind_{C \times P_0}^{C \rtimes P}(V)$, where $V \in \Irr_{\bC}(C \times P_0)$ corresponds to $W$, from which the claim follows as $C \times P_0$ is supersolvable, hence monomial by \cite[\S 8.5 Thm.\ 16]{SERRE}.

    Suppose instead that $W$ has S-type $\bR$, so $V \coloneqq W|_{C \times P_0}$ is irreducible. We adapt the proof of \cite[\S 8.5 Thm.\ 16]{SERRE}, and first show the claim when $C$ is trivial. We proceed by induction on the order of $P_0$. If $P_0$ is trivial then $P_0 \leq P$ is the Real group $C_1 \leq C_2$ and there is nothing to prove. Suppose instead that $P_0$ is non-trivial, so by induction we may assume that $W$, or equivalently $V$, is faithful.
    
    If $P_0$ is abelian, then because $V$ is irreducible, $W$ has dimension $1$, and we're done. If instead $P_0$ is not abelian, then $1 \neq P_0 / Z(P_0)$, and as $Z(P_0)$ is characteristic, $Z(P_0) \leq P$ is normal, so we can view $P_0 / Z(P_0) \leq P / Z(P_0)$. Because $P / Z(P_0)$ is a $2$-group, the normal subgroup $P_0 / Z(P_0)$ intersects $Z(P / Z(P_0))$ non-trivially, and so taking any element of order $2$ we can find a cyclic order $2$ normal subgroup $Q \triangleleft P / Z(P_0)$ which is contained in $P_0 / Z(P_0)$. Let $A$ denote the preimage of $Q$ in $P_0$. Because $A / Z(P_0)$ is cyclic, $A$ is abelian, and because $Q$ is normal in $P / Z(P_0)$, $A$ is normal in $P$.

    In summary, we have $A \leq P_0$, which is abelian, normal in $P$, and properly contains $Z(P_0)$. Consider $W|_A$. The representation $W|_{P_0}$ is faithful, so $W|_A$ cannot be isotypic: if it were then each $a \in A$ would have $\rho_V(a)$ a scalar matrix, and so commute with each $\rho_V(x)$ for $x \in P_0$, and so we would have $A \subset Z(P_0)$ by faithfulness, a contradiction. Therefore, the canonical decomposition $W|_A = V_1 \oplus \cdots \oplus V_k$ into isotypic factors necessarily has $k > 1$.
    
    Because $A$ is normal in $P$, there is a group action of $P$ on the set of isomorphism classes of representations of $A$: for $U \in \Rep_{\bC}(A)$ and $z \in P$, $z * U$ is $U$ as a $\bR$-vector space, with action of $a \in A$ by $z^{-1} a z$ and complex structure $\lambda * u \coloneqq \lambda u$ if $z \in P_0$ and $\lambda * u = \overline{\lambda} u$ is $z \in P \setminus P_0$. The action of each $z \in P$ defines an isomorphism $\rho_V(z) \colon z * W|_A \xrightarrow{\sim} W|_A$, which induces a permutation action of $P$ on the set $\{V_1, ... , V_k\}$. Let $S \coloneqq \Stab_P(V_1)$, noting that $S < P$ because $k > 1$ and the action is transitive ($W$ being irreducible). The natural map
    \[
        \bC_{\star \hspace{-0.1em}}[G] \otimes_{\bC_{\star \hspace{-0.1em}}[S]} V_1 \rightarrow W
    \]
    is surjective, the action being transitive, and therefore an isomorphism, both sides having the same dimension. In particular, $W \cong \Ind_S^G(V_1)$, and we are done by induction.

    This proves the claim when $C = 1$. In general, if $W$ has S-type $\bR$, then $W|_{C \times P_0}$ is irreducible, and because $C$ is abelian $C$ acts on $W|_{C \rtimes P_0}$ through some character $\chi \colon C \rightarrow \bC^{\times}$. In particular, $W|_{P_0}$ is still irreducible, and also of S-type $\bR$ (admitting the extension $W|_P$ to an antilinear representation). We can therefore write $W|_{P} = \Ind_S^P(\chi)$, for some Real subgroup $S \leq P$ and $1$-dimensional $U \in \Irr_{\bC}^{\hspace{-0.1em} \star }(S)$.
    
    Define an antilinear representation $\chi \otimes U$ of $C \rtimes S$, which is $U$ as an antilinear $S$-representation, with action of $C$ by scalar multiplication by $\chi$. The is well-defined: for $y \in C$ and $z \in S$, then if $z \in S_0$, $zy = y z$ in $C \rtimes S$ and indeed
    \[
        \rho_{\chi \otimes U}(z) \rho_{\chi \otimes U}(y)u
        = \rho_U(z)(\chi(y)(u)) = \chi(y) \rho_U(z)(u) =  \rho_{\chi \otimes U}(y) \rho_{\chi \otimes U}(z)u.
    \] 
    If on the other hand $z \in S \setminus S_0$, then $zy = y^{-1}z$ in $C \rtimes S$, and indeed
    \[
    \rho_{\chi \otimes U}(z) \rho_{\chi \otimes U}(y)u
        = \rho_U(z)(\chi(y)(u)) = \overline{\chi(y)} \rho_U(z)(u) =  \rho_{\chi \otimes U}(y^{-1}) \rho_{\chi \otimes U}(z)u.
    \]
    Then the fact that $W \cong \Ind_{C \rtimes S}^{C \rtimes P}(\chi \otimes U)$ follows from comparing characters.
\end{proof}

\begin{remark}
    Proposition \ref{prop:realelementarygrpsaremonomial} says that Real-elementary groups are \emph{Real-monomial}. Classically, supersolvable groups are monomial, and any monomial group is solvable, so monomial groups are somewhere between these two notions. In the antilinear setting, from computations it seems that a Real group $H \leq G$ is Real-monomial whenever $H$ is supersolvable, but we were unable to prove this.
\end{remark}

The second main result we need is an antilinear induction theorem for Real-elementary subgroups.

\begin{thm}\label{thm:antilinearinductionthm}
For any $C_2$-graded group $G$ and $\chi \in \Char^{\hspace{0.1em} \star}(G, \bC)$, there are $C_2$-graded subgroups $G_i$ of $G$, $n_i \in \bZ$, and $\chi_i \in \Char^{\hspace{0.1em} \star}(G_i, \bL)$ such that each $G_i$ is either ungraded or Real-elementary and
\[
    \chi = \sum_{i} n_i \cdot \Ind_{G_i}^G(\chi_i).
\]
\end{thm}

\subsection{Proof Of Theorem \ref{thm:antilinearinductionthm}}

In this section we prove Theorem \ref{thm:antilinearinductionthm}. The claim is trivial if $G$ is ungraded, so suppose that $G$ is non-trivially graded. Let $\sR^{\hspace{0.1em} \star}_{\bZ}(G)$ denote the Grothendieck group of antilinear representations of $G$. Let $I_1$ denote the image of the $\bZ$-linear map
\[
    \Ind_{\sI_1} \coloneqq \bigoplus_{H \leq \sI_1} \Ind_H^G \colon \bigoplus_{H \in \sI_1} \sR_{\bZ}(H) \rightarrow \sR^{\hspace{0.1em} \star}_{\bZ}(G)
\]
where $\sI_1$ is the collection of ungraded subgroups of $G$, and let $I_2$ denote the image of the $\bZ$-linear map
\[
    \Ind_{\sI_2} \coloneqq \bigoplus_{S \in \sI_2} \Ind_S^G \colon \bigoplus_{S \in \sI_2} \sR_{\bZ}^{\hspace{0.1em} \star}(S) \rightarrow \sR^{\hspace{0.1em} \star}_{\bZ}(G),
\]
where $\sI_2$ is the collection of Real-elementary subgroups of $G$. From the correspondence between irreducible representations of $G_0$ and antilinear representations of $G$ (Theorem \ref{thm:threeoptions}), we have that $I_1$ has finite index in $\sR^{\hspace{0.1em} \star}_{\bZ}(G)$ equal to $2^k$, where $k$ is the number of irreducible antilinear representations of $G$ of S-type $\bR$. Therefore, it is sufficient to show that $I \coloneqq I_1 + I_2$ has odd index in $\sR^{\hspace{0.1em} \star}_{\bZ}(G)$, for then $I$ must have both odd and even index, and thus $I = \sR^{\hspace{0.1em} \star}_{\bZ}(G)$. Because $I$ of $\sR^{\hspace{0.1em} \star}_{\bZ}(G)$ is an ideal by the projection formula, Lemma \ref{lem:projectionformula}, we are reduced to the following.

\begin{prop}\label{prop:imageoddindex}
If $|G| = 2^r m$ with $m$ odd, then $m \in I$.
\end{prop}

To prove Proposition \ref{prop:imageoddindex} we adapt the proof of \cite[\S 12.7]{SERRE} to our setting. Set $n \coloneqq |G_0|$, so $n = 2^{r-1} m$, and let $A$ denote the subring of $\overline{\bZ}$ generated by the $n$th roots of $1$. There is an induction map
\[
A \otimes_{\bZ} \Ind_{\sI} \colon \bigoplus_{S \in \sI} A \otimes_{\bZ} \sR_{\bZ}^{\hspace{0.1em} \star}(S) \rightarrow A \otimes_{\bZ} \sR^{\hspace{0.1em} \star}_{\bZ}(G),
\]
and just as in \cite[\S 10.2, Lem.\ 5]{SERRE} we have the following, from which we see that to prove Proposition \ref{prop:imageoddindex} it is sufficient to show that $m \in A \otimes_{\bZ} I$.

\begin{lemma}\label{lem:cyclotomicintegers}
   The image of $A \otimes_{\bZ} \Ind_{\sI}$ is $A \otimes_{\bZ} I$, and $(A \otimes I) \cap \sR^{\hspace{0.1em} \star}_{\bZ}(G) = I$.
\end{lemma}

Let $G_0 /\!/ G_0$ denote the set of conjugacy classes of $G_0$, and let $G_0 /\!/ G$ denote the set of Real conjugacy classes \cite[\S 1.3]{RUMTAY} of $G_0$ with respect to the Real structure $G$. Because antilinear representations of $G$ are Real class function and are determined by their restriction to $G_0$ the injective ring homomorphism
\[
\Phi_0 \colon A \otimes_{\bZ} \sR_{\bZ}(G_0) \hookrightarrow \Fun(G_0 /\!/ G_0, A),
\]
to the ring of functions from $G_0 /\!/ G_0$ to $A$ restricts to an injective ring homomorphism
\[
\Phi \colon A \otimes_{\bZ} \sR^{\hspace{0.1em} \star}_{\bZ}(G) \hookrightarrow \Fun(G_0 /\!/ G, A).
\]
to the ring of functions from $G_0 /\!/ G$ to $A$. The first ingredient in the proof of Proposition \ref{prop:imageoddindex} is the following sufficient criterion for such a function be in $A \otimes_{\bZ} I$.

\begin{lemma}\label{lem:intvaluedconstantfunctions}
    Suppose that $\varphi \in \Fun(G_0 /\!/ G, A)$ has values in $|G| \cdot \bZ$. Then $\varphi \in \Phi(A \otimes_{\bZ} I_1)$. 
\end{lemma}

In the proof we use the following action of $C_2 = \langle \epsilon \rangle$ on $\Fun(G_0 /\!/ G, A)$. Fix $\omega \in G \setminus G_0$, and set 
\[
(\epsilon * \varphi)(x) \coloneqq \varphi(\omega x^{-1} \omega^{-1}).
\]
With this definition, from \cite[Prop.\ 2.8]{RUMTAY} we have that for any $\varphi_0 \in A \otimes_{\bZ} \sR^{\hspace{0.1em} \star}_{\bZ}(G)$,
\[
\Phi(\Ind_{G_0}^G(\varphi_0)) = \Phi_0(\varphi_0) + \epsilon * \Phi_{0}(\varphi_0).
\]
\begin{proof}
    Let $\varphi_0 \coloneqq \varphi_0 / 2$, which is a Real class function with values in $n \bZ$. By \cite[\S10.3, Lem.\ 6]{SERRE}, $\varphi_0 = \Phi_0(x_0)$ for some $x_0 \in A \otimes_{\bZ} \sR_{\bZ}(G_0)$. Set $x \coloneqq \Ind_{G_0}^G(x_0) \in A \otimes_{\bZ} I_1$. We have that
    \begin{align*}
        \Phi(x) = \Phi(\Ind_{G_0}^G(x_0)) =  \Phi_0(x_0) + \epsilon * \Phi_{0}(x_0) = \varphi_0 + \epsilon * \varphi_0 = \varphi,
    \end{align*}
    because $\varphi_0$ is a Real class function, and thus fixed by the $C_2$ action.
\end{proof}

We now use $I_2$ and Real-elementary subgroups, which we haven't needed until now.

\begin{defn}
For any $x \in G_0$ of odd order, set
\[
Z(x) \coloneqq \{g \in G_0 \mid gxg^{-1} = x\} \sqcup \{z \in G \setminus G_0 \mid zxz^{-1} = x^{-1}\},
\]
fix a $2$-Sylow subgroup $P$ of $Z(x) / \langle x \rangle$, and let $R(x)$ be the preimage of $P$ in $Z(x)$.
\end{defn}

Given $x \in G_0$ of odd order, either $R(x)$ is ungraded (contained in $G_0$) or $R(x)$ is Real-elementary (Definition \ref{def:realelementary}), because $R(x)$ is the semi-direct product of $\langle x \rangle$ and any $2$-Sylow subgroup of $R(x)$.

\begin{lemma}\label{lem:localspecialelt}
    Suppose that $x \in G_0$ has odd order. Then there is $\varphi_x \in A \otimes_{\bZ} \sR^{\hspace{0.1em} \star}_{\bZ}(R(x))$ such that
    \[
    \varphi \coloneqq \Ind_{R(x)}^G(\varphi_x) \in A \otimes_{\bZ} I
    \]
    satisfies the following:
    \begin{enumerate}
        \item $\Phi(\varphi)$ is $\bZ$-valued,
        \item $\Phi(\varphi)(x) \not\equiv 0 \!\! \mod 2$,
        \item $\Phi(\varphi)(y) = 0$ for any odd order $y \in G_0$ which is not Real conjugate to $x$.
    \end{enumerate}
\end{lemma}

\begin{proof}
    We have that $R(x)_0$ is a $2$-elementary subgroup of $G_0$, and thus by \cite[\S 10.3, Lem.\ 8]{SERRE} there is $\varphi_{x,0} \in A \otimes_{\bZ} \sR_{\bZ}(R(x)_0)$ such that, setting $\varphi_0 \coloneqq \Ind_{R(x)_0}^{G_0}(\varphi_{x,0})$, 
    \begin{enumerate}
        \item $\Phi(\varphi_0)$ is $\bZ$-valued,
        \item $\Phi(\varphi_0)(x) \not\equiv 0 \!\! \mod 2$,
        \item $\Phi(\varphi_0)(y) = 0$ for any odd order $y \in G_0$ which is not conjugate to $x$.
    \end{enumerate}
    
    There are now two cases. Suppose first that $R(x)$ is ungraded, so $R(x) \subset G_0$ and $R(x) = R(x)_0$. Set $\varphi_x \coloneqq \varphi_{x,0}$ and $\varphi \coloneqq \Ind_{R(x)}^G(\varphi_x) = \Ind_{G_0}^{G}(\varphi_0)$. We have that
    \[
        \Phi(\varphi) = \Phi_0(\varphi_0) + \epsilon * \Phi_{0}(\varphi_0),
    \]
    and so $\Phi(\varphi)$ is $\bZ$-valued. To verify the second condition, suppose that $y \in G_0$ has odd order. Then
    \[
        \Phi(\varphi)(y) = \Phi_0(\varphi_0)(y) + \Phi_{0}(\varphi_0)(\omega y^{-1} \omega^{-1}).
    \]
    When $y = x$, $\omega x^{-1} \omega^{-1}$ is not $G_0$-conjugate to $x$ because $Z(x) \subset G_0$, so $\Phi_{0}(\varphi_0)(\omega x^{-1} \omega^{-1}) \equiv 0 \!\! \mod 2$ and $\Phi(\varphi)(x) \equiv 0 \!\! \mod 2$. When $y$ is not Real conjugate to $x$, then both $y$ and $\omega y^{-1} \omega^{-1}$ are not $G_0$-conjugate to $x$, hence $\Phi(\varphi)(y) = 0$.

    Suppose instead that $R(x)$ is a Real subgroup of $G$, and choose a $2$-Sylow subgroup $P$ of $R(x)$ so that we can view $R(x) = C \rtimes P$ with $R(x)_0 = C \times P_0$, where $P_0$ is an index two subgroup of $P$. The element $\varphi_{x,0} \in A \otimes_{\bZ} \sR_{\bZ}(R(x)_0)$ is constructed as
    \[
        \varphi_{x,0} = \sum_{\chi} \chi(x^{-1}) \otimes (\chi \otimes \chi_1)
    \]
    as $\chi$ runs over the irreducible characters of $C$, and $\chi_1$ is the trivial character of $P_0$. Because $\chi_1$ extends to an antilinear representation of $P$, the following the final part of the proof of Proposition \ref{prop:realelementarygrpsaremonomial}, each $\chi \otimes \chi_1$ is the character of an object of $\sR^{\hspace{0.1em} \star}_{\bZ}(R(x))$ and we can view $\varphi_{x,0}$ as the restriction of $\varphi_x \in A \otimes_{\bZ} \sR^{\hspace{0.1em} \star}_{\bZ}(R(x))$. Setting $\varphi \coloneqq \Ind_{R(x)}^G(\varphi_x)$ we have by Lemma \ref{lem:inductioncompatibility} that $\Phi(\varphi) = \Phi_0(\varphi_0)$, and therefore $\Phi(\varphi)$ is $\bZ$-valued, $\Phi(\varphi)(x) \not\equiv 0 \!\! \mod 2$, and $\Phi(\varphi)(y) = 0$ for any odd order $y \in G_0$ which is not Real conjugate to $x$, as the fact that $Z(x) \not\subset G_0$ means that the Real conjugacy of $x$ is the same as the $G_0$-conjugacy class of $x$.
\end{proof}

\begin{cor}\label{cor:particularelt}
There exists $\varphi \in A \otimes_{\bZ} I$ such that
\begin{enumerate}
    \item $\Phi(\varphi)$ is $\bZ$-valued,
    \item $\Phi(\varphi)(x) \not\equiv 0 \!\! \mod 2$ for any $x \in G_0$.
\end{enumerate}
\end{cor}

\begin{proof}
    Let $x_1, ... , x_k$ denote representations for the Real conjugacy classes of $G_0$ which consist of elements of odd order. Let $\varphi_i \in A \otimes_{\bZ} I$ be the element provided by Lemma \ref{lem:localspecialelt}, and set $\varphi \coloneqq \sum_i \varphi_i$. Clearly $\Phi(\varphi)$ is $\bZ$-valued, this being true of each $\Phi(\varphi_i)$. For the second condition, let $x \in G_0$, and write $x'$ for the $2$-regular part of $x$ as in \cite[\S 10.1]{SERRE}. From \cite[\S 10.3, Lem.\ 7]{SERRE}, $
    \Phi(\varphi_i)(x) \equiv \Phi(\varphi_i)(x') \!\! \mod 2$ for any $i = 1, ... , k$. Furthermore, $x'$ is Real conjugate among the $x_1, ... , x_k$ to exactly one element $x_i$, and therefore
    \[
        \Phi(\varphi)(x) \equiv \sum_j \Phi(\varphi_j)(x') \equiv  \Phi(\varphi_i)(x_i) \not\equiv 0 \!\! \mod 2. \qedhere
    \]
\end{proof}

We are now in a position to prove Proposition \ref{prop:imageoddindex}.

\begin{proof}[Proof of Proposition \ref{prop:imageoddindex}]
    Let $\varphi \in A \otimes_{\bZ} I_2$ be as in Corollary \ref{cor:particularelt}, and set $N \coloneqq 2^r$, the order of $(\bZ / 2^{r+1} \bZ)^{\times}$. Consider the function $\phi \coloneqq m \Phi(\varphi^N - 1) \in \Fun(G_0 / \! / G, A)$. Because $\Phi(\varphi)$ is $\bZ$-valued, the same is true for $\phi$, and because $\Phi(\varphi)(x) \not\equiv 0 \!\! \mod 2$ for any $x \in G_0$, $\Phi(\varphi)(x)^N \equiv 1 \!\! \mod 2$, hence $\phi$ has values in $|G| \cdot \bZ$. By Lemma \ref{lem:intvaluedconstantfunctions} there is $\psi \in A \otimes_{\bZ} I_1$ with $\Phi(\psi) = \phi$, and so by the injectivity of $\Phi$, $\psi = m\varphi^N - m$. We have that $m \varphi^N \in A \otimes_{\bZ} I$ because $A \otimes_{\bZ} I$ is an ideal of $A \otimes_{\bZ} \sR^{\hspace{0.1em} \star}_{\bZ}(G)$ and $\varphi \in A \otimes_{\bZ} I$ by Corollary \ref{cor:particularelt}. It then follows that $m = \psi - m \varphi^N \in A \otimes_{\bZ} I$, from which the claim follows from Lemma \ref{lem:cyclotomicintegers}.
\end{proof}

\subsection{Proof of Theorem \ref{thm:fundamentalrepsexplicit}}

We are now in a position to prove Theorem \ref{thm:fundamentalrepsexplicit}. We first show that $\sG_{\bC / \bR} \subset \sH$, where $\sH \coloneqq \sH^{1} \sqcup \sH^{-1}$, for $\sH^1$ the ungraded list
\[
\sH^{1} = \{C_n \mid n \geq 1\},
\]
and $\sH^{-1}$ the non-trivially graded list
\[
\sH^{-1} = \{C_n \leq D_{2n} \mid n \geq 1\}.
\]
We then compute exactly $\sG_{\bC / \bR} \subset \sH$ to be the sublist as described by Theorem \ref{thm:fundamentalrepsexplicit}.

We show that $\sG_{\bC / \bR} \subset \sH$ by induction on the order (see Definition \ref{def:orderC2gradedgroup}) of the $C_2$-graded group $G$. When the $o(G) = 1$, the only option for $G$ is that $G = C_1$ with the trivial grading, and this lies in both $\sG_{\bC / \bR}$ and $\sH$. Suppose that $\sG_{\bC / \bR, \leq n} \subset \sH_{\leq n}$ for some $n \geq 1$, and take $G \in \sG_{\bC / \bR, n+1}$. By Proposition \ref{prop:minimallistexistsuniquely}, there is some $\chi \in \Char^{\hspace{0.1em} \star}(G, \bC)$ such that you cannot write
    \[
    \chi = \sum_{i} n_i \cdot \Ind_{G_i}^G(\Inf_{Q_i}^{G_i}(\psi_i)),
    \]
where $n_i \in \bZ$, $G_i$ are $C_2$-graded subgroups of $G$, $G_i \twoheadrightarrow Q_i$ are $C_2$-graded quotients, $\psi_i \in \Char^{\hspace{0.1em} \star}(Q_i, \bC)$, and each $o(Q_i) < o(G)$.

Suppose that $n$ is even, so $G$ is ungraded. By Brauer's induction theorem we can write $\chi$ as a linear combination of characters induced from one-dimensional representations of ungraded subgroups, each of which, being one-dimensional, is inflated from a representation of a cyclic group. In particular, if $G$ were not cyclic, then it would fail to satisfy the above property, and so we must have that $G \in \sH^{1}$.

Suppose instead that $n$ is odd, so $G$ is non-trivially graded. By Theorem \ref{thm:antilinearinductionthm}, we immediate see that $G$ is Real-elementary. Decomposing $\chi$ into irreducible components, we may also assume that $\chi$ is the character of an irreducible antilinear representation of $G$. Then by Proposition \ref{prop:realelementarygrpsaremonomial}, $\chi$ is one-dimensional. As $G$ satisfies the above property, $\chi$ must be faithful, and therefore $G \cong (C_{m} \leq D_{2m})$ by Lemma \ref{lem:1dimfftypeRisdihedral}, where $m \coloneqq (n+1)/2$. In particular, $G \in \sH^{-1}$, and thus by induction $\sG_{\bC / \bR} \subset \sH$.

Now we consider exactly which elements of $\sH$ are in $\sG_{\bC / \bR}$. If $G \in \sH^1$, so $G = C_n$ for some $n \geq 1$, then all subquotients of $G$ are also ungraded, and thus $G \in \sG_{\bC / \bR}$ if and only if $n \equiv 2 \!\! \mod 4$ by Theorem \ref{thm:wheniscyclicnotneeded}. Suppose instead $G \in \sH^{-1}$, so $G = D_{2n}$ for some $n \geq 1$. From Lemma \ref{lem:dihedralrealisable} the characters of antilinear representations of $G$ are the same as the characters of the ungraded subgroup $C_n$. The Real structure $C_n \leq D_{2n}$ has the property that every ungraded subgroup $C_m \leq C_n$ lies in the Real subgroup $(C_m \leq D_{2m}) \leq (C_n \leq D_{2n})$, and any Real quotient of $(C_n \leq D_{2n})$ is of the form $(C_n \leq D_{2n}) \twoheadrightarrow (C_m \leq D_{2m})$. In particular, by Lemma \ref{lem:inductioncompatibility} the $\bZ$-submodule of $\Char^{\hspace{0.1em} \star}(D_{2n}, \bC)$ generated by characters of proper $C_2$-graded subquotients under inflation and induction is the same as the $\bZ$-submodule of $\Char(C_{n}, \bC)$ generated by proper subquotients of $C_n$ under inflation and induction. We therefore have that the answer is the same as the ungraded $C_n$ case, and so $(C_n \leq D_{2n}) \in \sG_{\bC / \bR}$ if and only if $n \equiv 2 \!\! \mod 4$ by Theorem \ref{thm:wheniscyclicnotneeded}.

\subsection{Proof of Theorem \ref{thm:antilinearcyclotomicfinal}}\label{sect:proofofcyclorealise}

From Theorem \ref{thm:fundamentalrepsexplicit} and Lemma \ref{lem:dihedralrealisable} there are $X_1, ... , X_n \in \Rep_{E}^{\hspace{0.1em} \star}(G)$, with characters $\psi_1, ... ,\psi_n$ such that
\[
\chi_W = \sum_{i = 1}^n a_{i} \psi_i
\]
for some $a_i \in \bZ$. If $W$ has S-type $\bC$ or $\bH$ we are done by Lemma \ref{lem:onlycareabouttypeR}. Otherwise, $W$ has S-type $\bR$, and therefore $\langle \chi_W, \chi_W \rangle = 1$, so
\[
1 = \langle \chi_W, \chi_W \rangle = \sum_{i = 1}^n a_{i} \langle \chi_W, \psi_i \rangle.
\]

Let $V \in \Irr_{\bC}(H)$ correspond to $W$, and let $\tilde{V} \in \Irr_E(H)$ be a model of $V$ over $E$, using \cite[\S 12.3]{SERRE}. Let $\tilde{W} \in \Irr_{E}^{\hspace{0.1em} \star}(G)$ correspond to $\tilde{V}$, and note that $\chi_W = \chi_V = \chi_{\tilde{V}}$. We have that $m_{E_+}^E(\tilde{V}) \in \{1,2\}$, and $m_{E_+}^E(\tilde{V}) \mid \langle \chi_{\tilde{V}}, \chi_U \rangle$ for any $U \in \Rep_{E}^{\hspace{0.1em} \star}(G)$ by \cite[Prop.\ 9.3(4)]{TAYSchur}. In particular, if $m_{E_+}^E(\tilde{V}) = 2$, then each term $\langle \chi_W, \psi_i \rangle$ is divisible by two, which is impossible as the terms sum to $1$. Therefore, $m_{E_+}^E(\tilde{V}) = 1$, so $\chi_{\tilde{W}} = \chi_{\tilde{V}}$, and $\tilde{W}_{\bC} \cong W$, as their characters are the same. \qedsymbol

\subsection{Recovery of Serre's Induction Theorem over $\bR$}\label{sect:recoveryofclassical}

Let $G$ be a $C_2$-graded group with trivial grading. Then every $C_2$-graded subgroup and $C_2$-graded quotient also has the trivial grading, and we see that Theorem \ref{thm:fundamentalrepsexplicit} recovers Brauer's Induction Theorem \cite[\S 10.5]{SERRE}, that any complex representation of a finite group is a $\bZ$-linear combination of representations induced from one-dimensional representations of its subgroups.

We now show how Theorem \ref{thm:fundamentalrepsexplicit} also recovers Serre's Induction Theorem \cite[Lem.\ 73.19]{CR2}, for $\bR$-linear representations.

\begin{cor}\label{cor:recoverserre}
    Let $H$ be a finite group, and $\chi \in \Char(H, \bR)$. Then
    \[
    \chi = \sum_{i} n_i \Ind_{H_i}^H(\psi_i),
    \]
    where $n_i \in \bZ$, $H_i \leq H$, and $\psi_i$ is inflated from an $\bR$-representation of a cyclic or dihedral group. 
\end{cor}

Before giving the proof, we need the following result regarding induction from subgroups of $H \times C_2$ which are not the product of subgroups of $H$ and $C_2$.

Recall that by Goursat's Lemma, subgroups of $H \times C_2$ come in three flavours: 
\begin{enumerate}
    \item $K \times 1$, for $K \leq H$,
    \item $K \times C_2$, for $K \leq H$,
    \item $S(K,u) \coloneqq \{(k, u(k)) \mid k \in K\}$, where $K \leq H$ and $u \colon K \twoheadrightarrow C_2$.
\end{enumerate}

It is these subgroups of the third class that do most of the work when Theorem \ref{thm:antilinearinductionthm} is applied to split Real structures. For these subgroups, note that there is an canonical isomorphism $K \xrightarrow{\sim} S(K,u)$, and using this we have a restriction functor
\[
    (-)|_{\bR[K]} \colon \Rep_{\bC}^{\hspace{0.1em} \star}(S(K,u)) \rightarrow \Rep_{\bR}(K).
\]

\begin{lemma}\label{lem:twistedsubgroupinductionsplit}
    The diagram
% https://q.uiver.app/#q=WzAsNCxbMCwwLCJcXFJlcF97XFxiQ31ee1xcaHNwYWNlezAuMWVtfSBcXHN0YXJ9KFMoSyx1KSkgIl0sWzEsMCwiXFxSZXBfe1xcYkN9XntcXGhzcGFjZXswLjFlbX0gXFxzdGFyfShIIFxcdGltZXMgQ18yKSJdLFswLDEsIlxcUmVwX3tcXGJSfShLKSJdLFsxLDEsIlxcUmVwX3tcXGJSfShIKSJdLFswLDIsIigtKXxfe1xcYlJbS119Il0sWzIsMywiXFxJbmRfS15IIiwyXSxbMSwzLCJcXHNpbSJdLFswLDEsIlxcSW5kX3tTKEssdSl9XntIIFxcdGltZXMgQ18yfSJdXQ==
\[\begin{tikzcd}[column sep=large]
	{\Rep_{\bC}^{\hspace{0.1em} \star}(S(K,u)) } & {\Rep_{\bC}^{\hspace{0.1em} \star}(H \times C_2)} \\
	{\Rep_{\bR}(K)} & {\Rep_{\bR}(H)}
	\arrow["{\Ind_{S(K,u)}^{H \times C_2}}", from=1-1, to=1-2]
	\arrow["{(-)|_{\bR[K]}}", from=1-1, to=2-1]
	\arrow["\sim", from=1-2, to=2-2]
	\arrow["{\Ind_K^H}"', from=2-1, to=2-2]
\end{tikzcd}\]
    commutes up to isomorphism.
\end{lemma}

\begin{remark}
    This isomorphism can probably be made natural, but the above is sufficient for our purposes.
\end{remark}

\begin{proof}
    Let $V \in \Rep_{\bC}^{\hspace{0.1em} \star}(S(K,u))$ and let $\chi$ denote the character of $V$. If $\psi$ denote the character of $V|_{\bR[K]}$, then by \cite[\S 2.6]{RUMTAY} $\psi(k) = \chi(k) + \overline{\chi}(k)$ when $k \in K_0$ and $\psi(k) = 0$ when $k \in K \setminus K_0$. We thus compute that $\Ind_K^H(\psi)$ has character with
    \begin{align*}
    |K| \cdot \Ind_K^H(\psi)(h) &= \sum_{\substack{x \in H \\ h \in xKx^{-1}}} \psi(x^{-1} h x), \\
    &= \sum_{\substack{x \in H \\ h \in xK_0x^{-1}}} \chi(x^{-1} h x) + \overline{\chi(x^{-1} h x)}, \\
    &= \sum_{\substack{x \in H \\ h \in xK_0x^{-1}}} \chi(x^{-1} h x) + \chi((x\omega)^{-1} h (x \omega)), \\
    &= 2 \cdot \sum_{\substack{x \in H \\ h \in xK_0x^{-1}}} \chi(x^{-1} h x),
    \end{align*}
    and therefore
    \[
    \Ind_K^H(\psi)(h) = \frac{1}{|K_0|} \sum_{\substack{x \in H \\ h \in xK_0x^{-1}}} \chi(x^{-1} h x),
    \]
    which is exactly the character of $\Ind_{S(K,u)}^{H \times C_2}(V)$ by Lemma \ref{lem:inductioncompatibility}. Then we see that
    \[
    \Ind_K^H(V|_{\bR[G]}) \cong \Ind_{S(K,u)}^{H \times C_2}(V)^{C_2},
    \]
    because they have the same character by the above and \cite[Lem.\ 6.2]{TAYSchur}. 
\end{proof}

\begin{proof}[Proof of Corollary \ref{cor:recoverserre}]
    Let $\chi \in \Char(H, \bR)$, and let $U \in \Rep_{\bR}(H)$ be the corresponding representation. We have that $\bC \otimes_{\bR} U \in \Rep_{\bC}^{\hspace{0.1em} \star}(H \times C_2)$, which also has character $\chi$ by \cite[Lem.\ 6.2]{TAYSchur}. From Theorem \ref{thm:fundamentalrepsexplicit} there are
    \begin{enumerate}
        \item ungraded subgroups $H_i \leq H$ and Real subgroups $S_j \leq H \times C_2$,
        \item $C_2$-graded quotients $H_i \twoheadrightarrow C_{n_i}$ and $S_j \twoheadrightarrow D_{2n_j}$,
        \item $\chi_i \in \Char(C_{n_i},\bC)$ and $\chi_j \in \Char^{\hspace{0.1em} \star}(D_{2n_j},\bC)$
    \end{enumerate}
    such that
    \[
        \chi = \sum_i m_i \Ind_{H_i \times 1}^{H \times C_2}(\chi_i) + \sum_j m_j \Ind_{S_j}^{H \times C_2}(\chi_j)
    \]
    for some $m_i, m_j \in \bZ$. For any $H_i$, there is a commutative diagram
% https://q.uiver.app/#q=WzAsOCxbMSwwLCJcXFJlcF97XFxiQ30oSF9pKSJdLFsyLDEsIlxcUmVwX3tcXGJDfShIKSJdLFszLDAsIlxcUmVwX3tcXGJDfV57XFxoc3BhY2V7MC4xZW19IFxcc3Rhcn0oSCBcXHRpbWVzIENfMikiXSxbMywyLCJcXFJlcF97XFxiUn0oSCkiXSxbMSwyLCJcXFJlcF97XFxiUn0oSF9pKSJdLFs0LDEsIlxcUmVwX3tcXGJDfShIKSJdLFswLDAsIlxcUmVwX3tcXGJDfShDX3tuX2l9KSJdLFswLDIsIlxcUmVwX3tcXGJSfShDX3tuX2l9KSJdLFswLDIsIlxcSW5kX3tIX2kgXFx0aW1lcyAxfV57SCBcXHRpbWVzIENfMn0iXSxbMCw0LCIoLSl8X3tcXGJSW0hfaV19IiwyXSxbNCwzLCJcXEluZF9LXkgiLDJdLFswLDEsIlxcSW5kX3tIX2l9XntIfSIsMl0sWzEsMiwiXFxJbmRfe0ggXFx0aW1lcyAxfV57SCBcXHRpbWVzIENfMn0iLDJdLFsyLDMsIlxcc2ltIiwyXSxbMSwzLCIoLSl8X3tcXGJSfSIsMl0sWzMsMiwiKC0pXntDXzJ9IiwyLHsic3R5bGUiOnsiYm9keSI6eyJuYW1lIjoibm9uZSJ9LCJoZWFkIjp7Im5hbWUiOiJub25lIn19fV0sWzIsNSwiKC0pfF9IIl0sWzMsNSwiKC0pX3tcXGJDfSIsMl0sWzYsMCwiXFxJbmZfe0Nfbn1ee0hfaX0iXSxbNiw3LCIoLSl8X3tcXGJSW0Nfe25faX1dfSIsMl0sWzcsNCwiXFxJbmZfe0Nfe25faX19XntIX2l9IiwyXV0=
\[\begin{tikzcd}
	{\Rep_{\bC}(C_{n_i})} & {\Rep_{\bC}(H_i)} && {\Rep_{\bC}^{\hspace{0.1em} \star}(H \times C_2)} \\
	&& {\Rep_{\bC}(H)} && {\Rep_{\bC}(H)} \\
	{\Rep_{\bR}(C_{n_i})} & {\Rep_{\bR}(H_i)} && {\Rep_{\bR}(H)}
	\arrow["{\Inf_{C_n}^{H_i}}", from=1-1, to=1-2]
	\arrow["{(-)|_{\bR[C_{n_i}]}}"', from=1-1, to=3-1]
	\arrow["{\Ind_{H_i \times 1}^{H \times C_2}}", from=1-2, to=1-4]
	\arrow["{\Ind_{H_i}^{H}}"', from=1-2, to=2-3]
	\arrow["{(-)|_{\bR[H_i]}}"', from=1-2, to=3-2]
	\arrow["{(-)|_H}", from=1-4, to=2-5]
	\arrow["\sim"', from=1-4, to=3-4]
	\arrow["{\Ind_{H \times 1}^{H \times C_2}}"', from=2-3, to=1-4]
	\arrow["{(-)|_{\bR}}"', from=2-3, to=3-4]
	\arrow["{\Inf_{C_{n_i}}^{H_i}}"', from=3-1, to=3-2]
	\arrow["{\Ind_K^H}"', from=3-2, to=3-4]
	\arrow["{(-)^{C_2}}"', draw=none, from=3-4, to=1-4]
	\arrow["{(-)_{\bC}}"', from=3-4, to=2-5]
\end{tikzcd}\]
The top triangle commutes by the transitivity of induction, and the two triangles on the right by \cite[Lem.\ 6.2]{TAYSchur} and \cite[Lem.\ 6.3]{TAYSchur}. The left-hand square and bottom-left triangle commute because restriction from $\bC$ to $\bR$ commutes with induction and inflation. In particular, each character $\Ind_{H_i \times 1}^{H \times C_2}(\chi_i)$, obtained by going along the top from $\Rep_{\bC}(C_{n_i})$ to $\Rep_{\bC}(H)$, is the same as the character obtained by going around the other outside route, and therefore in the image of of $\Rep_{\bR}(C_{n_i})$.

For the remaining subgroups, note that if $S \leq H \times C_2$ is a Real subgroup which surjects onto $D_{2n}$ for some $n \geq 3$, then $S = S(K,u)$ for some $K \leq H$ and $u \colon K \twoheadrightarrow C_2$. Indeed, the only other option is that $S = K \times C_2$, and such a $q \colon K \times C_2 \twoheadrightarrow D_{2n}$ being $C_2$-graded would induce an isomorphism $K / \ker(q) \times C_2 \cong D_{2n}$, which is obviously impossible. We cannot have $n = 2$ by Theorem \ref{thm:fundamentalrepsexplicit}. If $n = 1$, then $S = 1 \times C_2$, and by Lemma \ref{lem:inductioncompatibility} any induced representation is a sum of copies of $\bC[H]$, which corresponds to the representation $\bR[H]$ of $H$ over $\bR$. However, this comes from inducing the trivial representation of the trivial subgroup of $H$ over $\bR$, and so is permitted. We may therefore restrict to $S_i = S(K_i,u_i)$, for which there is a commutative diagram,
% https://q.uiver.app/#q=WzAsNyxbMSwwLCJcXFJlcF97XFxiQ31ee1xcaHNwYWNlezAuMWVtfSBcXHN0YXJ9KFMoS19pLHVfaSkpICJdLFsyLDAsIlxcUmVwX3tcXGJDfV57XFxoc3BhY2V7MC4xZW19IFxcc3Rhcn0oSCBcXHRpbWVzIENfMikiXSxbMSwxLCJcXFJlcF97XFxiUn0oS19pKSJdLFsyLDEsIlxcUmVwX3tcXGJSfShIKSJdLFswLDAsIlxcUmVwX3tcXGJDfV57XFxoc3BhY2V7MC4xZW19IFxcc3Rhcn0oRF97Mm5fan0pICJdLFswLDEsIlxcUmVwX3tcXGJSfShEX3sybl9qfSkgIl0sWzMsMSwiXFxSZXBfe1xcYkN9KEgpIl0sWzAsMiwiKC0pfF97XFxiUltLX2ldfSJdLFsyLDMsIlxcSW5kX3tLX2l9XkgiLDJdLFsxLDMsIlxcc2ltIl0sWzAsMSwiXFxJbmRfe1MoS19pLHVfaSl9XntIIFxcdGltZXMgQ18yfSJdLFs0LDAsIlxcSW5mX3tEX3syIG5fan19XntTKEtfaSwgdV9pKX0iXSxbNCw1LCIoLSl8X3tcXGJSW0RfezJuX2p9XX0iLDJdLFs1LDIsIlxcSW5mX3tEX3sybl9qfX1ee0tfaX0iLDJdLFszLDYsIigtKV97XFxiQ30iLDJdLFsxLDYsIigtKXxfSCJdXQ==
\[\begin{tikzcd}[column sep=large]
	{\Rep_{\bC}^{\hspace{0.1em} \star}(D_{2n_j}) } & {\Rep_{\bC}^{\hspace{0.1em} \star}(S(K_i,u_i)) } & {\Rep_{\bC}^{\hspace{0.1em} \star}(H \times C_2)} \\
	{\Rep_{\bR}(D_{2n_j}) } & {\Rep_{\bR}(K_i)} & {\Rep_{\bR}(H)} & {\Rep_{\bC}(H)}
	\arrow["{\Inf_{D_{2 n_j}}^{S(K_i, u_i)}}", from=1-1, to=1-2]
	\arrow["{(-)|_{\bR[D_{2n_j}]}}"', from=1-1, to=2-1]
	\arrow["{\Ind_{S(K_i,u_i)}^{H \times C_2}}", from=1-2, to=1-3]
	\arrow["{(-)|_{\bR[K_i]}}", from=1-2, to=2-2]
	\arrow["\sim", from=1-3, to=2-3]
	\arrow["{(-)|_H}", from=1-3, to=2-4]
	\arrow["{\Inf_{D_{2n_j}}^{K_i}}"', from=2-1, to=2-2]
	\arrow["{\Ind_{K_i}^H}"', from=2-2, to=2-3]
	\arrow["{(-)_{\bC}}"', from=2-3, to=2-4]
\end{tikzcd}\]
the left-hand square commuting by definition, the right-hand square by Lemma \ref{lem:twistedsubgroupinductionsplit}, and the right-hand triangle by \cite[Lem.\ 6.2]{TAYSchur}.
\end{proof}

\subsection{Classification of Types Over Cyclotomic Fields}

We can use Theorem \ref{thm:antilinearcyclotomicfinal} to give a complete description of the possible types of representations over cyclotomic fields and to describe the type in terms of easily computable invariants.

In addition to the cyclotomic field $E = \bQ(\zeta_e)$ and its index two maximal totally real subfield $E_+$, we set $f \coloneqq \exp(G)$ when $\exp(H) > 2$ and $f \coloneqq 4$ when $\exp(H) \leq 2$, and similarly define $F \coloneqq \bQ(\zeta_f)$ and $F_+ \coloneqq F \cap \bR = \bQ(\zeta_f + \zeta_f^{-1})$. Because $H$ has index two in $G$, either $e = f$ or $2e = f$, and so either $(F,F_+) = (E,E_+)$ or $F$ and $F_+$ are degree two extensions of $E$ and $E_+$ respectively.

\begin{thm}\label{thm:15waysforcyclotomic}
Let $V \in \Irr_{E}(H)$ and let $Z \in \Irr_{\bC}(G)$ correspond to $V_{\bC}$. Then the types of $V$ for the extensions $E_+ \leq E$, $F_+ \leq F$ and $\bR \leq \bC$ are in one of the cases of Table \ref{tab:15typesQe}. \emph{
\begin{table}[h!] % remove [h!] if you don't want float
\centering
\bgroup
\caption{The Possible Types of $V$ over $\bQ(\zeta_e)$ and $\bQ(\zeta_f)$}
\label{tab:15typesQe}
\def\arraystretch{1.1}
\begin{tabular}{ c  c  c  c }
    \toprule
    $\bR \leq \bC$ & $\bQ(\chi_Z) \subset E$ & $E_+ \leq E$ & $F_+ \leq F$ \\
    \midrule 
 I & Yes & (1) & (1) \\
 I & No & (8) & (1) \\ \midrule
 II & Yes & (5) & (5) \\
 II & No & (9) & (5) \\ \midrule
 III & Yes & (2) & (2) \\ \midrule
 IV & Yes & (10) & (10) \\
 IV & No & (14) & (10) \\ \midrule
 V & Yes & (12) & (12) \\ \midrule 
 VI & Yes & (15) & (15) \\ \midrule 
 VII & Yes & (11) & (11) \\ \midrule
 VIII & Yes & (16) & (16) \\
 VIII & No & (24) & (16) \\ \midrule 
 IX & Yes & (18) & (18) \\
 IX & No & (22) & (18) \\ \midrule 
 X & Yes & (17) & (17) \\ \bottomrule 
\end{tabular}
\egroup
\end{table}
}
\end{thm}

Here $\bQ(\chi_Z)$ is the extension of $\bQ$ generated by the character values of $\chi_Z$.

\begin{remark}\label{rem:FSforcyclo}
    The type of $V_{\bC}$ is determined by three computable quantities: the classical and antilinear Frobenius-Schur indicators, and whether or not $\chi_V \cong {}^{\omega}\chi_V$. From the Theorem, the type of $V_{F}$ and $V_{\bC}$ are the same (just using different labels), and therefore the type of $V_{F}$ is also determined by these quantities. We also see that the type of $V$ over $E$ is determined by these quantities and whether or not $\bQ(\chi_Z) \subset E$. These are all computationally implementable, and we magma code for this at \cite{GITHUB}.
\end{remark}

\begin{remark}
    All cases of Theorem \ref{thm:15waysforcyclotomic} do actually occur! Indeed, the examples of Section \ref{sect:realise24ways} for these 15 cases are over $E_+ \leq E$.
\end{remark}

\begin{proof}
    The type of $V$ for the extension $\bL[H] \leq \bL_{\star \hspace{-0.1em}}[G]$ is the same for all three choices of $\bK \leq \bL$ by Theorem \ref{thm:antilinearcyclotomicfinal}. Similarly, the type of $V$ for the extension $\bK[H] \leq \bL[H]$ is the same for all three choices of $\bK \leq \bL$ by Theorem \ref{thm:classicalcyclotomic} (which itself follows from Theorem \ref{thm:antilinearcyclotomicfinal}, see Section \ref{sect:recoveryofclassical}). In particular, the type of $V_{\bC}$ is uniquely determined by the type of $V$ over $E$, as it is uniquely by the type of $V_{\bC}$ for the extensions $\bR[H] \leq \bC[H]$ and $\bC[H] \leq \bC_{\star \hspace{-0.1em}}[G]$ and whether or not $\overline{\chi}_V \cong {}^{\omega}\chi_V$. Furthermore, the types of $V_{F}$ for the extensions $F_+[G] \leq F[G]$ and $F[H] \leq F[G]$ are the same as the types of $V_{\bC}$ for the extension $\bR[G] \leq \bC[G]$ and $\bR[H] \leq \bC[G]$ respectively, by Theorem \ref{thm:classicalcyclotomic} and the fact that $F$ is a splitting field for $H$ and $G$. Therefore, the type of $V_{F}$ is equal to the type of $V_{\bC}$, any row of Table \ref{tab:24types} being determined by its first four columns.

    Note also that whether or not $\bQ(\chi_Z) \not\subset E$ is determined by the type of $V$ over $E$, as this is equivalent to $\End_{E[G]}(Z_0)$ being $2$-dimensional over $E$, where $Z_0 \in \Irr_{E}(H)$ corresponds to $V$. Indeed, if $e=f$, then neither can happen, and if $f = 2e$ then both are equivalent to being in Case B of Table \ref{tab:thewaysLHSsquare} for the extension $E[G] \leq F[G]$.
    
    Therefore, it suffices to show that the only types that can occur for $V$ over $E$ are those 15 listed in Table \ref{tab:15typesQe} out of the possible 24 of Theorem \ref{thm:24ways}. Cases (3), (6), (7), (13), (20) and (21) are ruled out immediately using Table \ref{tab:24types6} because $E$ is a splitting field for $H$, and so we are left with types (4), (19) and (23).

    If $V$ had type (4), then $V_{F}$ would have type $(1)$. If $Z_e \in \Irr_{E}(G)$ corresponds to $V$, then as $V$ has type (4), $Z_e$ is absolutely irreducible by Table \ref{tab:24types6}, and so $Z_{e, F}$ corresponds to $V_{F}$. Then $\Stab_{K_4}(Z_{e, F}) = \Stab_{K_4}(Z_e) = \langle c \rangle$ by Table \ref{tab:24types5} as $V$ has type (4), which contradicts the fact that $\Stab_{K_4}(Z_{e, F}) = \langle a \rangle$ because $V_{F}$ has type $(1)$, again by Table \ref{tab:24types5}.

    If $V$ had type (19), then $V_{F}$ would have type $(16)$. Similarly, if $Z_e \in \Irr_{E}(G)$ corresponds to $V$, then as $V$ has type (4), $Z_e$ is absolutely irreducible by Table \ref{tab:24types6}, and therefore $Z_{e,F}$ corresponds to $V_{F}$. We again obtain a contradiction: from Table \ref{tab:24types5} this would force $\Stab_{K_4}(Z_{e,F})$ to be simultaneously $\langle a \rangle$ and $\langle c \rangle$.

    Finally, if $V$ had type (23), then $V_{F}$ would have type $(16)$. If $U_e \in \Irr_{E_+}(G)$ corresponds to $V$, then as $V$ has type (23) $U_{e, E}$ is irreducible, and $\End_{E[G]}(U_{e, E})$ is $2$-dimensional over $E$. In particular, as $F$ splits $G$ then from Table \ref{tab:thewaysLHSsquare} of Section \ref{sect:classicalindextwofieldextoptions}
    \[
    (U_{e, E})_{F} = U_{e, F} = Z_1 \oplus Z_2
    \]
    for some non-isomorphic $Z_1, Z_2 \in \Irr_{F}(G)$. On the other hand, if $U_f \in \Irr_{F_+}(G)$ corresponds to $V_{F}$, then as $V_{F}$ has type $(16)$, $U_{f, F} = Z_3 \oplus Z_3$ for some $Z_3 \in \Irr_{F}(G)$ by Table \ref{tab:24types2}. However the irreducible factors of $U_{e, F_+}$ correspond to $V_{F}$, and therefore $U_f \hookrightarrow U_{e, F_+}$. In particular, 
    \[
        Z_3 \oplus Z_3 = U_{f, F} \hookrightarrow U_{e, F} = Z_1 \oplus Z_2,
    \]
    so $Z_1 \cong Z_3 \cong Z_2$, which contradicts the fact that $Z_1$ and $Z_2$ are non-isomorphic.
\end{proof}

\subsection{Applications to Classical Schur Indices}

We can give some application of Theorem \ref{thm:15waysforcyclotomic} to classical Schur indices for groups which admit an index two subgroup. In what follows, whenever $f > e$, let $\sigma$ denote the non-trivial element of $\Gal(F / E)$, which we also use to denote its restriction to the non-trivial element of $\Gal(F_+ / E_+)$.

\begin{prop}\label{prop:indextwosubgroupclassicalCextension}
    Suppose that $G$ is a finite group with an index two subgroup $H$ of exponent~$e$. For $Z \in \Irr_{\bC}(G)$, let $Z_f \in \Irr_{F}(G)$ and $Z_e \in \Irr_{E}(G)$ correspond to $Z$. Then $Z$ is in one of the two cases of Table \ref{tab:2optionsZQef}
    \emph{
\begin{table}[h!] % remove [h!] if you don't want float
\centering
\bgroup
\caption{The Possible Models for $Z$ over $\bQ(\zeta_e)$ and $\bQ(\zeta_f)$}
\label{tab:2optionsZQef}
\def\arraystretch{1.1}
\begin{tabular}{ c c c}
    \toprule 
    $\bQ(\chi_Z) \subset E$ & $Z_{f, \bC}$ & $Z_{e, F}$ \\ \midrule
 Yes & $Z$ & $Z_f$ \\
 No & $Z$ & $Z_f \oplus Z_f^{\sigma}$ \\ \bottomrule
\end{tabular}
\egroup
\end{table}
}
In particular $\overline{m}_{E}(Z) = 1$: $Z$ admits a model over $E(\chi_Z)$.
\end{prop}

\begin{proof}
    First, because $F$ splits $G$ \cite[\S 12.3]{SERRE}, $Z_{f, \bC} \cong Z$. If $E = F$ then we are in the first case, so suppose not, meaning that $2e = f$ and the extension $F / E$ has degree $2$.
    
    From Section \ref{sect:classicalindextwofieldextoptions}, there are three cases possible for the extension $E[G] \leq F[G]$, listed in Table \ref{tab:thewaysLHSsquare}. Case A corresponds to the first line of Table \ref{tab:2optionsZQef}, and Case B to the second, and therefore the claim is that Case C cannot occur. Indeed, in Case C, $\End_{E[G]}(Z_e)$ has dimension $4$ over $E$, from Table \ref{tab:thewaysLHSsquare} and the fact that $\End_{F[G]}(Z_f)$ has dimension $1$ over $F$. However this is impossible: $Z_e$ is also in one of the four cases of Proposition \ref{prop:algindressummary} for the extension $E[H] \subset E[G]$ (cf.\ Section \ref{sect:indextwosubgroup}), and hence $\End_{E[G]}(Z_f)$ has dimension $1$ or $2$, because $E$ splits $H$.
\end{proof}

The proof of Proposition \ref{prop:indextwosubgroupclassicalCextension} is relatively straightforward using Proposition \ref{prop:algindressummary}, and doesn't use Theorem \ref{thm:antilinearcyclotomicfinal} and its consequences. It could instead be proven using the classification of Theorem \ref{thm:15waysforcyclotomic} as we do now in the following, where the use of Theorem \ref{thm:15waysforcyclotomic} is now essential.

\begin{prop}\label{prop:indextwosubgroupclassicalRextension}
    Suppose that $G$ is a finite group with an index two subgroup $H$ of exponent~$e$. For $U \in \Irr_{\bR}(G)$, let $U_f \in \Irr_{F_+}(G)$ and $U_e \in \Irr_{E_+}(G)$ correspond to $U$. Then $U$ is in one of the two cases of Table \ref{tab:2optionsUQef}
    \emph{
\begin{table}[h!] % remove [h!] if you don't want float
\centering
\bgroup
\caption{The Possible Models for $U$ over $\bQ(\zeta_e)_+$ and $\bQ(\zeta_f)_+$}
\label{tab:2optionsUQef}
\def\arraystretch{1.1}
\begin{tabular}{c c c }
    \toprule
    $\bQ(\chi_U) \subset E_+$ & $U_{f, \bR}$ & $U_{e, F_+}$ \\ \midrule 
 Yes & $U$ & $U_f$ \\
 No & $U$ & $U_f \oplus U_f^{\sigma}$ \\ \bottomrule
\end{tabular}
\egroup
\end{table}
}
In particular $\overline{m}_{E_+}(U) = 1$: $U$ admits a model over $E_+(\chi_U)$.
\end{prop}

\begin{proof}
    By Theorem \ref{thm:classicalcyclotomic} we have that $U_{f, \bC} \cong U$. If $E = F$ then we are in the first case, so suppose not, meaning that $2e = f$ and the extension $F / E$ has degree $2$. From Section \ref{sect:classicalindextwofieldextoptions}, there are three possibilities for the extension $E_+[G] \leq F_+[G]$. Case A corresponds to the first line of Table \ref{tab:2optionsUQef}, and Case B to the second, and therefore the claim is that Case C cannot occur.
    
    To see this, let $V \in \Irr_{E}(H)$ correspond to $U_e$, and so $V_{F}$ corresponds to $U_f$. One can check using Table \ref{tab:24types6} that in each of the 15 cases of Table \ref{tab:15typesQe} the ratio
    \[
    \dim_{F} \End_{F[G]}(U_f) : \dim_{E} \End_{E[G]}(U_e)
    \]
    is either $1 : 1$ or $1 : 2$, and therefore that Case C, when the ratio would be $1 : 4$, never occurs.
\end{proof}

\begin{remark}
    Propositions \ref{prop:indextwosubgroupclassicalCextension} and \ref{prop:indextwosubgroupclassicalRextension} are not true without the existence of the index two subgroup. For example, if one instead assumes that $G$ has even order and takes $m \coloneqq \text{Exponent}(G) / 2$, then the conclusion of Theorem \ref{prop:indextwosubgroupclassicalRextension} is false. An example is given by SmallGroup(48,41). However, it seems that the statement remains is also true if one instead uses the assumption that $G$ has even order and takes $m \coloneqq |G|/2$ (we have checked this computationally up to $G$ of order 200).
\end{remark}

\section{Unitary Representations}\label{sect:unitaryreps}

In this section we no longer assume that $G$ is finite, but instead restrict our attention to unitary (antilinear) representations. This gives us several additional tools that make the theory reminiscent of that of a finite group.

\begin{defn}
  Let $(V,\tau)$ be a finite-dimensional hermitian vector space (over $\bC$). We write $\UU(V)$ for the unitary group of $V$,
  and $\UU^\ast (V)$ for the set of antilinear transformations $f: V\to V$, ``flipping'' the hermitian form, i.e.,
\[ \tau (f(x), f(y)) = \tau (y,x) = \overline{\tau (x,y)} \, . \]
  We write $\UR (V)\coloneqq \UU(V) \cup \UU^\ast (V)$, which we view as a $C_2$-graded group (cf.\ Section \ref{sect:fundreps}) under composition.

  An \emph{(antilinear) unitary representation} of a $C_2$-graded group $(G,\epsilon_G)$ is a homomorphism of $C_2$-graded groups $G\to \UR (V)$, where $V$ is a finite-dimensional hermitian vector space.
\end{defn}

Observe that an antilinear map will not preserve a hermitian form in general. Indeed, writing vectors as columns in an orthonormal basis, the form becomes $\tau (x,y) = x^T \overline{y}$, while an antilinear transformation is given by $x \mapsto A \overline{x}$. Then the flipping/preserving conditions become
\[ \tau (A\overline{x},A \overline{y}) = \overline{x}^T A^T \overline{A} y  \stackrel{?}{=}
\begin{cases}
	\tau(x,y) =  x^T \overline{y} & \ \mbox{preserving},\\
	\overline{\tau(x,y)} =  \overline{x}^T {y} & \ \mbox{flipping}.
\end{cases}\]
for all $x,y$. Clearly, the preserving condition leads to a contradiction by choosing various $x$ and $y$, while the flipping condition is
just $A^T \overline{A}=I$ which, for instance, all real orthogonal matrices satisfy.

Note that if $G$ is finite, all antilinear representations over $\bC/\bR$ are unitary by Weyl's Unitary Trick:
pick any positive definite hermitian form $\tau$ on $V$, then integrate it over $G$ to get a new form
\[ \widetilde{\tau} (x,y) = \sum_{g\in H} \tau (gx,gy) + \sum_{g\in G\setminus H} \tau (gy,gx)\, .\]

\subsection{Bohr compactification}%%{Bohr completion of a group}
The Bohr compactification of a group $G$ is the universal homomorphism $G\to \widehat{G}$ to a compact group.
In the literature it is often used for topological groups but here we are working only with discrete groups. 

The best way to think of the Bohr compactification $\widehat{G}$ is that it is the group one gets from Tannaka-Krein Duality, applied to the category of unitary (linear) representations of $G$. Let us adopt a more down-to-earth definition.

\begin{defn}\label{def:bohrcompactification}
The Bohr compactification $\widehat{G}$ is the closure of the image of $G$ in $\prod_V \UU(V)$ where the product is taken over isomorphism classed over irreducible unitary (linear) representations of $G$.
\end{defn}

From the definition, we have a continuous map $\alpha : G \to \widehat{G}$. Since the Bohr compactification is not really our focus, we will just list its standard properties. In what follows, we let $\beta \colon G \rightarrow \Prof(G)$ denote the profinite completion of $G$.
We write $\URep(\widehat{G})$ for the set of isomorphism classes of finite-dimensional continuous unitary representations of $\widehat{G}$, and $\UIrr(\widehat{G})$ for the subset of irreducible ones. We write the same for $G$, noting that the continuity requirement is vacuous and can be dropped.

\begin{prop}\label{prop:bohrcompactificationproperties} The following properties hold for a discrete group $G$.
  \begin{enumerate}
  \item $\alpha (G)$ is dense in $\widehat{G}$.
  \item If $G$ is abelian, then $\widehat{G}$ is abelian.
  \item If $\beta: G \to K$ is a homomorphism to a compact Hausdorff group $K$, then there is a unique continuous homomorphism $f: \widehat{G}\to K$ such that $\beta = f \circ \alpha$.
  \item In particular, the association $G \mapsto \widehat{G}$ is functorial, and is left adjoint to the forgetful functor from compact groups to discrete groups.
  \item Sending $(V, \rho) \mapsto (V, \rho \circ \alpha)$ defines bijections
  \[
  \URep(\widehat{G}) \xlongrightarrow{\sim} \URep(G), \qquad \UIrr(\widehat{G}) \xlongrightarrow{\sim} \UIrr(G).
  \]
  \item The natural map $\widehat{G} \rightarrow \Prof(G)$ induced by $\beta \colon G \rightarrow \Prof(G)$ induces an isomorphism 
  \[
  \pi_0(\widehat{G}) = \widehat{G} / \widehat{G}_0 \xrightarrow{\sim} \Prof(G),
  \]
  as topological groups, where $G_0$ is the connected component of the identity in $\widehat{G}$.
  \item If $H$ is a finite-index normal subgroup of $G$, then the induced map $\widehat{H} \rightarrow \widehat{G}$ is injective, the image $\widehat{H}$ is a finite-index normal subgroup of $\widehat{G}$, and $\alpha_G \colon G \rightarrow \widehat{G}$ induces an isomorphism
  \[
  G/H \cong \widehat{G}/ \widehat{H}.
  \]
  \end{enumerate}  
\end{prop}  

\begin{proof}
    The first two points follow from the definition, the second because when $G$ is abelian each factor $U(V)$ is isomorphic to $S^1$. Points (3) to (5) follow from \cite[Prop.\ 5]{BEK1}. The sixth point is \cite[Prop.\ 7]{BEK1}, and the final point follows from \cite[Prop.\ 9]{BEK1} and \cite[Prop.\ 10]{BEK1}.
\end{proof}
 
Note that the map $\alpha : G \to \widehat{G}$ is not necessarily injective. For example, if $G$ is a non-linear simple group, then $\widehat{G}$ is the trivial group. If $\alpha$ is injective, then the topology on $G$ induced from $\widehat{G}$ is discrete if and only if $G$ is finite because $\widehat{G}$ is compact Hausdorff and $G$ is dense in $\widehat{G}$. In particular, the Bohr compactification is not a topological compactification in the usual sense, and perhaps ``Bohr completion'' would be a more appropriate term here as well.

We also note the similarity between the Bohr compactification of a topological group $G$ and the Bohr completion of a $\bK$-algebra $A$ as defined in Section \ref{Bohr_completion} (for example, compare Definition \ref{def:bohrcompactification} with Lemma \ref{Bohr2}). Certain properties that hold for one may not translate obviously to the other: for example the correspondence of simple modules of Proposition \ref{prop:bohrcompactificationproperties} fails for Bohr completions of algebras (see Remark \ref{rem:bohrcompletionexample}), but one instead has Lemma \ref{lem:idemirrcorresp}.

\subsection{Frobenius-Schur indicators for compact groups} \label{section:FS_UIrr} 

If $G$ is a finite $C_2$-graded group, we have three Frobenius-Schur indicators in play \cite{RUMTAY, RUMTAY3}. The indicators involve summations over $G$ or $H$, the index 2 subgroup. For infinite groups, they work in the same way but only for unitary representations.

First, let $G$ be a compact Hausdorff topological group, and $H$ an index two (normal) subgroup of $G$. Let $\mu_G$, $\mu_H$ denote Haar measures on $G$ and $H$ normalised in such a way that $\mu_G(G) = 1$ and $\mu_H(H) = 1$. Define functions 
\begin{equation} \label{FS_formulas}
\sF,{\sF^{\sharp}} : \UIrr (H) \rightarrow \bC , \quad
\sF (V) =
\int_{H} \chi_V (x^2) \mu_G(dx), \quad
    {\sF^{\sharp}} (V) =
    \int_{G\setminus H} \chi_V (x^2) \mu_G(dx),
    \end{equation}

We note that for $V \in\UIrr (H)$, $V$ is finite-dimensional, so we may consider $V$ as an irreducible representation of $H$ viewed as an abstract group, and associate to $V$ a type $\fT(V)$, a D-type $\fD(V)$, an S-type $\fS(V)$ and a C-type $\fC(V)$.

The following is a straightforward generalisation to compact Hausdorff groups of the relationship between the Frobenius-Schur indicator and the classical and semilinear types.

\begin{prop}\label{prop:FSexttocompact}
    Suppose that $V \in \UIrr (H)$. Then 
    \[
        \sF(V) = \left\{
		\begin{array}{cl}
			1 & \mbox{if } \: \fC(V) = \textbf{\emph{A}} \\
            0 & \mbox{if } \: \fC(V) = \textbf{\emph{B}} \\
			-1 & \mbox{if } \: \fC(V) = \textbf{\emph{C}}
		\end{array}
        \right. ,
        \qquad 
        \sF^{\sharp}(V) = \left\{
		\begin{array}{cl}
			1 & \mbox{if } \: \fS(V) = \textbf{\emph{A}} \\
            0 & \mbox{if } \: \fS(V) = \textbf{\emph{B}} \\
			-1 & \mbox{if } \: \fS(V) = \textbf{\emph{C}}
		\end{array}
        \right. .
    \]
\end{prop}

\begin{proof}
    For $\sF(V)$, this is well-known, for example, see \cite[Prop.\ 6.8]{DieckBrocker}. For $\sF^{\sharp}(V)$, this follows exactly as in the proof of \cite[Thm.\ 4.2]{RUMTAY}, by computing $\sF^{\sharp}(V)$ for each possible D-type.
\end{proof}

\subsection{Frobenius-Schur indicators for discrete groups}

Now suppose that $G$ is a discrete, potentially infinite, group, with an index two subgroup $H$.

\begin{lemma}
    Let $V \in \UIrr(H)$, and let $\widehat{V} \in \UIrr(\widehat{H})$ be the corresponding irreducible unitary representation of $\widehat{H}$. Then $\fC(V) = \fC(\widehat{V})$ and $\fS(V) = \fS(\widehat{V})$.
\end{lemma}

\begin{proof}
    To ease notation, we write $(-)^{\wedge}$ instead $\widehat{(-)}$. We first show that $\fC(V) = \fC(V^{\wedge})$. Let $X^{\wedge} \in \Irr_{\bR}(\widehat{H})$ denote the irreducible $\bR$-linear representation corresponding to $V^{\wedge}$, which is continuous because $X^{\wedge}$ is a subrepresentation of the restriction to $\bR$ of $V^{\wedge} \oplus \overline{V^{\wedge}}$. We can pull back along $\alpha$ to obtain an $\bR$-linear representation $X$ of $H$, which is irreducible by continuity, because $H$ is dense in $\widehat{H}$. By construction, $X$ is a constituent of $V|_{\bR}$, as $X^{\wedge}$ is a constituent of $V^{\wedge}|_{\bR}$. Again by continuity, the restriction map
    \[
    \End_{\bR \widehat{H}}(X^{\wedge}) \rightarrow \End_{\bR H}(X)
    \]
    is an isomorphism, and thus $\fC(V) = \fC(V^{\wedge})$. A similar argument shows that $\fS(V) = \fS(V^{\wedge})$.
\end{proof}

Therefore, we can compute the D-type of $V$ using integration on $\widehat{H}$ and $\widehat{G}$, by Proposition \ref{prop:FSexttocompact}.
We note that this is not clear how to do this for a general group, where a useful description of the Bohr compactification is often not available.

\subsection{Frobenius-Schur indicators for amenable groups}
Countable amenable groups are interesting because
the integrals~\ref{FS_formulas} can be computed intrinsically, without a reference to the Bohr compactification.
As amenable groups are not our focus, we take the following definition, it being the most convenient (see \cite[Prop.\ 4.7.5]{CSCoornaert}). 

\begin{defn}
A discrete countable group $G$ is amenable if it admits a F{\o}lner sequence $(F_n)_n$. The latter consists of a sequence of non-empty finite subsets $F_n\subseteq G$ for $n \geq 1$ such that
\[ \forall g\in G \quad \lim_{n \rightarrow \infty} \frac{|gF_n \oplus F_n|}{|F_n|} = 0\, ,\]
where $A \oplus B$ is the symmetric difference of the subsets.
\end{defn}

Let $\sM(G) \subset \ell^{\infty}(G,\mathbb{R})^\ast$ denote the space of means on $G$, consisting of linear functionals $\mu \in \ell^{\infty}(G,\mathbb{R})^\ast$ which are positive ($\mu(f) \geq 0$ whenever $f \geq 0$) and satisfy $\mu(1)=1$.
It follows from the Banach–Alaoglu Theorem that $\sM(G)$ is compact in the weak-$\ast$ topology \cite[Thm.\ 4.2.1]{CSCoornaert}.
The F{\o}lner sequence gives a sequence of averaging functionals
\[
  \mu_n \in \sM(G), \qquad
  \mu_n(f) = \frac{1}{|F_n|}\sum_{x\in F_n} f(x),
\]
which has a cluster point $\mu_G \in \sM(G)$ in the weak-$\ast$ topology by the compactness of $\sM(G)$.
We note that $\mu_G$ is not unique in general. For $\mu \in \ell^{\infty}(G,\mathbb{R})^*$ and $f \in \ell^{\infty}(G,\mathbb{C})$, we define 
\[
\mu(f) \coloneqq \mu(\Re(f)) + i \mu(\Im(f))
\]
and denote
\[
\int_G f(x) d \mu(x) \coloneqq \mu(f).
\]

Recall that $\alpha\colon G\to\widehat G$ denotes the Bohr compactification of $G$,
and $\mu_{\widehat G}$ denotes the normalized Haar measure on $\widehat G$.
We are going to use the following lemma only for matrix elements.
We suspect it is known to specialists but we could not a reference, so we provide a full proof.

\begin{lemma}\label{lem:folnerintegral}
Let $G$ be a discrete countable amenable group with F{\o}lner sequence
$(F_n)_n$, $(\mu_n)_n$ -- the corresponding sequence of means,
$\mu_G$ -- any weak-$\ast$ cluster point of $(\mu_n)_n$. Then for every
$f\in C(\widehat G,\mathbb{C})$, we have
\[
  \int_{\widehat G} f(z)\,d\mu_{\widehat G}(z)
  \;=\; \int_G (f\circ\alpha)(x)\,d\mu_G(x)
  \;=\; \lim_{n\to\infty}\mu_n(f\circ\alpha).
\]
\end{lemma}

\begin{proof}
  Intuitively, it suffices to prove the statement for matrix elements because they are dense (Stone-Weierstrass Theorem), and it is true
  for matrix elements because averaging defines an orthogonal projection $L^2 (\widehat{G},\bC) \to L^2 (\widehat{G},\bC)^G = \bC$.
  Now let us give full details.

Pick  $(V,\pi) \in \UIrr(\widehat{G})$, so that $(V,\rho \coloneqq \pi\alpha) \in \UIrr({G})$. Set
\[
  A_n=\frac{1}{|F_n|}\sum_{x\in F_n}\rho(x),
  \qquad
  P=\int_{\widehat G}\pi(z)\,d\mu_{\widehat G}(z)
  \qquad\in\mathrm{End}_{\bC}(V) \, .
  \]
We claim that the sequences of matrices $A_n$ converges to the matrix $P$.

To prove the claim, examine $P$. It is an endomorphism of $V$ as a $\widehat{G}$-module. By Schur's Lemma $P=\alpha I$.
Moreover, $P(v)$ is $\widehat{G}$-invariant: $\pi (g) (P(v)) = P(v)$. Hence, $P = 0$ if $V$ is non-trivial, and $P=I$ if $V$ is trivial.


To examine $A_n$, we use the operator norm on $\mathrm{End}_{\bC}(V)$.
Each $\rho(x)$ is unitary, so $\| \rho(x) \| =1$ and $\|A_n\|\leq 1$.
Thus, all $A_n$ belong to the (compact) unit ball and the sequence $(A_n)$ has convergent subsequences.
To prove the claim, it remains to show that $P$ is the limit of every convergent subsequence.

Since for each
$g\in G$ we have $\rho(g)\sum_{x\in F_n}\rho(x)=\sum_{y\in gF_n}\rho(y)$, we can estimate 
\[
  \|\rho(g)A_n-A_n\|
  =\frac{1}{|F_n|}\Big\|\sum_{y\in gF_n}\rho(y)-\sum_{x\in F_n}\rho(x)\Big\|
  \leq\frac{|gF_n\oplus F_n|}{|F_n|}\xrightarrow[n\to\infty]{}0 .
  \]
  Thus, the limit of any convergent subsequence is a linear operator $V \to V^G = V^{\widetilde{G}}$.
  The claim follows immediately for a non-trivial $V$. For the trivial $V$ we need an additional observation that $A_n = I$.

  At this stage, we have proved that
\[
  \int_{\widehat G} f(z)\,d\mu_{\widehat G}(z)
  \;=\; \lim_{n\to\infty}\mu_n(f\circ\alpha), 
\]
  if $f$ is a matrix coefficient of a finite-dimensional unitary representation, i.e., %%the functions $f$ of the form
  \[ f (x)=\langle\pi(x)v,w\rangle, \ \mbox{ for some } v,w\in (V,\pi) \in \URep(\widehat{G}) \, , \]
  since we can extend the statement of the proposition from matrix coefficients of irreducible representations to their finite linear combinations.
  Let us extend this to all continuous functions.
  
Let $\mathcal T\subset C(\widehat G,\mathbb{C})$ be the linear span of the
matrix coefficients of all finite-dimensional continuous unitary
representations of $\widehat G$.
By the combination of Peter-Weyl and Stone-Weierstrass theorems,
$\mathcal T$ is dense in $C(\widehat G,\mathbb{C})$ in the sup norm.
Now consider
$f\in C(\widehat G,\mathbb{C})$.
Given  $\varepsilon>0$, choose $h\in\mathcal T$
such that  $\|f-h\|_\infty<\varepsilon / 3$.
By the triangle inequality, 
\[
  \Big|\mu_n(f\circ\alpha)-\int_{\widehat{G}} f\,d\mu_{\widehat G}\Big|
  \leq
  \Big|\mu_n((f-h)\circ\alpha)\Big|
  +\Big|\mu_n(h\circ\alpha)-\int_{\widehat{G}} h\,d\mu_{\widehat G}\Big|
  +\Big|\int_{\widehat{G}}(h-f)\,d\mu_{\widehat G}\Big|. 
  \]
Since
$\|(f-h)\circ\alpha\|_\infty\leq\|f-h\|_\infty < \varepsilon/3$  
and $\int\int_{\widehat{G}} d \mu_{\widehat G} =1$, we can estimate the first and third terms
\[
  \Big|\mu_n(f\circ\alpha)-\int_{\widehat{G}} f\,d\mu_{\widehat G}\Big|
  < \frac{2}{3} \varepsilon
   +\Big|\mu_n(h\circ\alpha)-\int_{\widehat{G}} h\,d\mu_{\widehat G}\Big|.
  \]
  Now choose $N$ such that for $n>N$ we have $|\mu_n(h\circ\alpha)-\int_{\widehat{G}} h\,d\mu_{\widehat G}| < \varepsilon/3$.
  Then for $n>N$,
  \[|\mu_n(f\circ\alpha)-\int_{\widehat{G}} f\,d\mu_{\widehat G}| < \varepsilon \, . \]

  It remains to examine the middle integral $\int_G (f\circ\alpha)(x)\,d\mu_G(x)$. Obviously, our functions belong to
$\ell^\infty(G,\mathbb{R})$ and for any
$h\in\ell^\infty(G,\mathbb{R})$ such that $\mu_n(h)\to c$ for some number $c$, we have
$\mu_G(h)=c$.
\end{proof}

If $G$ is an abstract group, and $H$ is an index two subgroup, then $G$ is amenable if and only if $H$ is amenable. Given an odd element $\omega \in G \setminus H$ and a F{\o}lner sequence $F_n$ of $H$, then $(F_n \cup \omega F_n)_n$ is a F{\o}lner sequence for $G$.
 
\begin{prop}\label{prop:FSamenable}
    Suppose that $H$ is a discrete countable amenable group, $G$ is a Real structure on $H$, $F_n$ is a F{\o}lner sequence for $H$ and $\omega \in G \setminus H$. Then for any $V \in \UIrr(H)$ the quantities
    \[
    \sF_{\Sigma}(V) \coloneqq \lim_{n \rightarrow \infty} \left( \frac{1}{|F_n|} \sum_{x \in F_n} \chi_V(x^2)\right),
    \qquad
    \sF_{\Sigma}^{\sharp}(V) \coloneqq \lim_{n \rightarrow \infty} \left( \frac{1}{|F_n|} \sum_{x \in F_n} \chi_V((\omega x)^2)\right)
    \]
    converge and satisfy
    \[
        \sF_{\Sigma}(V) = \left\{
		\begin{array}{cl}
			1 & \mbox{if } \: \fC(V) = \textbf{\emph{A}} \\
            0 & \mbox{if } \: \fC(V) = \textbf{\emph{B}} \\
			-1 & \mbox{if } \: \fC(V) = \textbf{\emph{C}}
		\end{array}
        \right. ,
        \qquad 
        \sF_{\Sigma}^{\sharp}(V) = \left\{
		\begin{array}{cl}
			1 & \mbox{if } \: \fS(V) = \textbf{\emph{A}} \\
            0 & \mbox{if } \: \fS(V) = \textbf{\emph{B}} \\
			-1 & \mbox{if } \: \fS(V) = \textbf{\emph{C}}
		\end{array}
        \right. .
    \]
\end{prop}

\begin{proof}
    Let $\widehat{V} \in \UIrr(\widehat{G})$ be the unique extension of $V$ from Proposition \ref{prop:bohrcompactificationproperties}, and let $f \in C(\widehat{G}, \bC)$ defined by $f(x) \coloneqq f(x) \coloneqq \chi_{\widehat{V}}(x^2)$, so that 
    \[
    \sF(\widehat{V}) = \int_{\widehat{H}} f(x) d \mu_{\widehat{H}}(x), \qquad \sF^{\sharp}(\widehat{V}) = \int_{\widehat{G} \setminus \widehat{H}} f(x) d \mu_{\widehat{G}}(x).
    \]
    By Proposition \ref{prop:FSexttocompact} it suffices to show that $\sF(\widehat{V}) = \sF_{\Sigma}(V)$ and $\sF^{\sharp}(\widehat{V}) = \sF^{\sharp}_{\Sigma}(V)$. By definition,
    \[
    \sF_{\Sigma}(V) = \lim_{n\to\infty}\mu_{H,n}(f\circ\alpha_H),
    \]
    and therefore that $\sF(\widehat{V}) = \sF_{\Sigma}(V)$ follows directly from Lemma \ref{lem:folnerintegral}. For the second equality, note that by Proposition \ref{prop:bohrcompactificationproperties}(7), the induced map $\widehat{H} \rightarrow \widehat{G}$ is injective and we can view $\widehat{H}$ as an index two normal subgroup of $\widehat{G}$. Then we again use Lemma \ref{lem:folnerintegral} for both $H$ and $G$ to compute that
    \begin{align*}
        \sF^{\sharp}_{\Sigma}(V)
        &= \lim_{n\to\infty}\mu_{G,n}(f\circ\alpha_G) - \lim_{n\to\infty}\mu_{H,n}(f\circ\alpha_H), \\
        &= \int_{\widehat{G}} f(x) d \mu_{\widehat{G}}(x) - \int_{\widehat{H}} f(x) d \mu_{\widehat{H}}(x), \\
        &= \int_{\widehat{G} \setminus \widehat{H}} f(x) d \mu_{\widehat{G}}(x), \\
        &= \sF^{\sharp}(\widehat{V})
    \end{align*}
    as required.
\end{proof}

\section{Dependence of the D-Type on the Real Structure}\label{sect:deponrealstr}

In this section we consider how, for a group $H$, the D-type of some $V \in \Irr_{\bC}(H)$ depends on a given Real structure $G$.

\begin{lemma}\label{lem:indepK4action}
    The action of $K_4$ on $\Irr_{\bC}(H)$ depends only on the coupling of $G$.
\end{lemma}

\begin{proof}
    This is clear, as this action is defined in terms of conjugation by an element of $G \setminus H$.
\end{proof}

\begin{defn}
    For any $V \in \Irr_{\bC}(H)$, we write $\mathfrak{D}(V; G)$ for the D-type of $V$ with respect to the real structure $G$. When $H$ is finite, we set
    \[
    \sF(V ; G) = \frac{1}{|H|} \sum_{g \in G \setminus H} \chi_V(g^2).
    \]
\end{defn}

\begin{cor}
    Let $V \in \Irr_{\bC}(H)$. Then the properties that
    \begin{itemize}
        \item $\mathfrak{D}(V; G) = \text{\emph{III}}$,
        \item $\mathfrak{D}(V; G) = \text{\emph{IV}}$,
        \item $\mathfrak{D}(V; G) = \text{\emph{VII}}$,
        \item $\mathfrak{D}(V; G) = \text{\emph{X}}$,
    \end{itemize}
    depend only on the coupling of the Real structure $G$.
\end{cor}

\begin{proof}
    By Lemma \ref{lem:indepK4action}, $\Stab_{K_4}(V)$ depends only on the coupling. Furthermore, the C-type of $V$ being $\bR$, $\bC$ or $\bH$ depends only on $G$, and not on the Real structure at all. The D-types which are determined by these two pieces of data are those listed above, by Table 2 of \cite{RUMTAY3}.
\end{proof}

Recall that by Lemma \ref{lem:elementary1} any Real structure on $H$ is of the form $G = G(\alpha, h)$, where $\alpha \in \Aut(H)$ and $h \in H$ with $\alpha^2 = \Ad(h)$ and $\alpha(h) = h$. Let us fix such a model. Then a set of representatives of all Real structures on $G$ with the same coupling $\theta \coloneqq \overline{\alpha}$ as $G(\alpha, h)$ is $G(\alpha, zh)$, where $z$ varies $\HH^2(Z_{\theta})$ which we view as the quotient $HH^2(Z_{\theta}) \coloneqq \{z \in Z(H) \mid \alpha(z) = z\} / \{z\alpha(z) \mid z \in Z(H)\}$.

\begin{prop}\label{prop:variationofFS}
    Suppose that $H$ is finite and $V \in \Irr_{\bC}(H)$ with character $\chi$. Then
    \[
        \sF(V; G(\alpha, zh)) = \frac{\chi(z)}{\dim(V)} \cdot \sF(V; G(\alpha, h)).
    \]
    In particular, if $\mathfrak{D}(V; G(\alpha, h))$ is not equal to \emph{III, IV, VII} or \emph{X}, $\chi(z) \in \{\pm \dim(V) \}$. In this case, if $\chi(z) = \dim(V)$, then 
    \[
    \mathfrak{D}(V; G(\alpha, h)) = \mathfrak{D}(V; G(\alpha, zh)),
    \]
    and if $\chi(z) = -\dim(V)$, then $\{\mathfrak{D}(V; G(\alpha, h)), \mathfrak{D}(V; G(\alpha, zh))\}$ is one of $\{\text{\emph{I}}, \text{\emph{II}}\}$, $\{\text{\emph{V}}, \text{\emph{VI}}\}$, $\{\text{\emph{VIII}}, \text{\emph{IX}}\}$.
\end{prop}

\begin{proof}
    We can compute
    \begin{align*}
         |G| \cdot \sF(V; G(\alpha, h)) &= \sum_{g \in G(\alpha, zh) \setminus G} \chi(g^2), \\
         &= \sum_{g \in G} \chi((gz)^2), \\
         &= \sum_{g \in G} \chi(g \alpha(g) h).
    \end{align*}
    Replacing $h$ by $zh$, we then have that
    \begin{align*}
        |G| \cdot \sF(V; G(\alpha, zh)) &= \sum_{g \in G} \chi(g \alpha(g) h z), \\
        &= \frac{\chi(z)}{\dim(V)} \sum_{g \in G} \chi(g \alpha(g) h),
    \end{align*}
    using the fact that $\dim(V) \chi(gy) = \chi(g) \chi(y)$ for any central element $y$ of $G$, because $V$ is irreducible and so $y$ acts by a scalar. The remaining points then follow from Table 2 of \cite{RUMTAY3}.
\end{proof}

\begin{remark}
    Suppose that $\mathfrak{D}(V; G(\alpha, h))$ is not equal to III, IV, VII or X, and let $z \in Z(H)$. Because $[G(\alpha, h)] = [G(\alpha, z\alpha(z)h)]$, $\sF(V; G(\alpha, zh)) = \sF(V; G(\alpha, h))$ and therefore by the above we must have that $\chi(z \alpha(z)) = \dim(V)$, or equivalently that $\chi(z) \chi(\alpha(z)) = \dim(V)^2$. One can also see directly that this is true, because outside of D-types III, IV, VII or X, $\overline{\chi} = \chi(\alpha(-))$ because $c \in K_4$ stabilises $V$, hence $\chi(z) \chi(\alpha(z)) = |\chi(z)|^2 = \dim(V)^2$. 

    We also know that $\chi(z) = \pm \dim(V)$, by Proposition \ref{prop:variationofFS}, and therefore $\chi(z) = \chi(\alpha(z))$.
\end{remark}

This allows us to say identify immediately many $V \in \Irr_{\bC}(H)$ that will have S-type $\bC$.

\begin{cor}
    Suppose that $H$ is finite, $V \in \Irr_{\bC}(H)$, and $G$ is a Real structure on $H$ with coupling $\theta$. Then then if $\chi(z) \not\in \bR$ for some $z \in \HH^2(Z_{\theta})$, or $\chi(z) \neq \chi(\alpha(z))$ for some $z \in Z(H)$ and lift $\alpha$ of $\theta$, then $\sF(V, \widehat{G}) = 0$.
\end{cor}
